\documentclass[11pt,letterpaper]{article}

\usepackage[margin=1in]{geometry}
\usepackage{lmodern}
\usepackage[T1]{fontenc}
\usepackage[utf8]{inputenc}
\usepackage{microtype}
\usepackage{amsmath,amssymb,amsthm,mathtools}
\usepackage{graphicx}
\usepackage{booktabs}
\usepackage{tabularx}
\usepackage{caption}
\usepackage{xcolor}
\usepackage{fancyhdr}
\usepackage{hyperref}

\definecolor{jedcblue}{RGB}{24, 62, 114}
\hypersetup{
  colorlinks=true,
  linkcolor=jedcblue,
  citecolor=jedcblue,
  urlcolor=jedcblue,
  pdftitle={Observationally Equivalent Algorithms: What Prices Reveal About Trader Populations, and What Supervisors Need},
  pdfauthor={Pratik Niroula, Prashant Niroula, Sujan Khatiwada}
}

\setlength{\parskip}{4pt plus 1pt minus 1pt}
\setlength{\parindent}{1.2em}
\linespread{1.12}

\pagestyle{fancy}
\fancyhf{}
\fancyhead[L]{\small\scshape Niroula, Niroula, and Khatiwada}
\fancyhead[R]{\small\scshape Observationally Equivalent Algorithms}
\fancyfoot[C]{\small\thepage}
\renewcommand{\headrulewidth}{0.4pt}

\fancypagestyle{plain}{
  \fancyhf{}
  \fancyfoot[C]{\small\thepage}
  \renewcommand{\headrulewidth}{0pt}
}

\newtheorem{theorem}{Theorem}[section]
\newtheorem{proposition}[theorem]{Proposition}
\newtheorem{corollary}[theorem]{Corollary}
\newtheorem{lemma}[theorem]{Lemma}
\theoremstyle{definition}
\newtheorem{remark}[theorem]{Remark}
\newtheorem{definition}[theorem]{Definition}

\newcommand{\E}{\mathbb{E}}
\newcommand{\R}{\mathbb{R}}
\newcommand{\KL}{\mathrm{KL}}
\newcommand{\TV}{\mathrm{TV}}
\newcommand{\Var}{\mathrm{Var}}
\newcommand{\supp}{\mathrm{supp}}

\begin{document}

\thispagestyle{empty}
\begin{center}
  {\LARGE\bfseries Observationally Equivalent Algorithms:\\[4pt]
  What Prices Reveal About Trader Populations,\\[4pt]
  and What Supervisors Need\par}
  \vspace{1.2em}
  {\large Pratik Niroula\footnote{Department of Computer Information Science, Minnesota State University, Mankato, MN 56001, USA. Email: \texttt{pratik.niroula@mnsu.edu}.} \qquad
  Prashant Niroula\footnote{Faculty of Management, Tribhuvan University, Kirtipur, Nepal. Email: \texttt{prashantniroula20@gmail.com}.} \qquad
  Sujan Khatiwada\footnote{Department of Electrical, Computer and Robotics Engineering, Minnesota State University, Mankato, MN 56001, USA. Email: \texttt{sujan.khatiwada@mnsu.edu}.}}
  \vspace{0.8em}

  {\normalsize November 2026}
\end{center}

\vspace{0.6em}
\begin{abstract}
\noindent
Algorithmic feedback traders read prices and, by trading, move them. We show that a heterogeneous trader population affects public prices, order flows, and dealer inventories solely through its aggregate feedback kernel, both under mechanical impact and in a stationary inventory-dealer equilibrium. Separating $K$ trader groups whose memory rates lie within $\Delta$ requires of order $\Delta^{-(4K-2)}$ price observations—a barrier that persists arbitrarily close to instability. Consequently, a uniform cut in feedback can be calibrated directly from prices and, without class tags, is the unique minimax supervisory policy for any convex loss and strictly convex, label-symmetric cost (weakly dominating every price-based policy under weak convexity). Targeted restrictions require class-tagged flow reports, which eliminate the spacing penalty for weight recovery. However, when restricted traders can route orders through alternative labels, class tags lose all value as relabeling becomes cheap, while untagged surveillance induces strategic crowding that drives equilibrium supervisory information to zero at rate $N^{-2/(p-2)}$ as the record $N$ grows. Only owner-linked audit trails escape both incentive failures. A calibration to Taiwan Stock Exchange tagged institutional flows confirms that prices pin down a uniform cut tightly while leaving targeted halts unidentified.
\end{abstract}

\vspace{0.5em}
\noindent\textbf{Keywords:} algorithmic trading; feedback trading; aggregate response; observational equivalence; market supervision; disclosure.\\
\noindent\textbf{JEL Classification:} C32, D83, G14, G18.\\
\noindent\textbf{Online Supplement and Replication:} Technical lemmas, extended proofs, numerical robustness checks, and Lean~4 machine-checked formalizations of deterministic algebraic steps (with probabilistic and asymptotic steps proved on paper) are provided in the Online Supplement and public replication archive (\url{https://github.com/Pratikn03/Fianince}). Extensions to breach cascades, strategic insiders, and nonlinear or learning traders appear in a companion note (Niroula et al., 2026b).

\newpage
\setcounter{page}{1}

\section{Introduction}

\textbf{The economic question.} A financial market supervisor may know with high precision how strongly a market reacts to past prices while knowing almost nothing about which traders generate that reaction. These are fundamentally different objects of inference. If a policy scales every trader group's feedback response by the same factor, the aggregate market response is a sufficient statistic to evaluate the intervention. If instead a supervisor seeks to halt or restrict a specific trader class, that class's individual contribution is essential. This paper characterizes the exact boundary between what public prices reveal about trader populations and what requires class-tagged or owner-linked supervisory reporting.

\textbf{Main results.} We embed heterogeneous feedback-trading algorithms into both a reduced-form linear market and a microfounded inventory-dealer equilibrium. Under a common public signal and fixed causal rules, the trader population enters all public market outcomes—transaction prices, total order flows, insider profits, and dealer inventories—solely through its aggregate flow map (Theorem~\ref{thm:sufficiency}). For linear rules, this map is the \emph{feedback kernel}, the moment sequence of the population's merged memory-rate measure. In Section~\ref{sec:equilibrium}, we derive price impact and pricing-error persistence endogenously from competitive, risk-averse overlapping-generations dealers who bear inventory risk. Equilibrium prices follow the exact same closed-loop recursion with effective impact and persistence and unchanged memory rates (Theorems~\ref{thm:inv-eq} and~\ref{thm:eq-id}). Thus, observational equivalence and the sample-complexity barriers below are equilibrium phenomena rather than artifacts of mechanical price impact.

Five economic results organize the analysis:
\begin{enumerate}
  \item \emph{Weak identification of composition.} Closely spaced trader types are severely under-identified from prices. Two moment-matched $K$-type populations whose memory rates lie within $\Delta$ produce aggregate-flow differences of order $\Delta^{2K-1}$, yielding a Gaussian information lower bound of order $\Delta^{-(4K-2)}$ observations—whether the supervisor observes prices passively or conducts active impulse probes (Theorem~\ref{thm:cost}). This barrier extends to arbitrary smooth causal rule families (Theorem~\ref{thm:smooth}), matches a Prony and minimum-distance upper bound on separated shrinking clusters (Theorem~\ref{thm:recovery}), and does not vanish as the market approaches instability (Figure~\ref{fig:near-instability}).
  \item \emph{Separation between market-wide and compositional questions.} Market-wide linear functionals of the merged measure (such as discounted feedback or single-lag taps) can be estimated at the standard parametric rate $\sigma^2 \epsilon^{-2}$ independent of memory spacing, whereas compositional weights require $\Delta^{-(4K-2)}$ observations (Theorem~\ref{thm:cheap} and Corollary~\ref{cor:sep}). In a Gaussian feedback market, the Kullback--Leibler divergence between two candidate populations equals half the cumulative excess mean-squared forecast error from confusing them (Theorem~\ref{thm:forecast-kl}); hence composition is unnecessary for optimal price forecasting but indispensable for targeted attribution.
  \item \emph{Minimax optimality of uniform policies without tags.} Because prices cannot distinguish label permutations, a uniform cut across all feedback activity weakly dominates every price-based policy against the worst labeling consistent with the record for any convex loss and label-symmetric cost—and is the unique minimax supervisory policy whenever the cost is strictly convex (Theorem~\ref{thm:minimax}). With verified partial tags partitioning classes into groups, a policy that is uniform within each tagged group is minimax-optimal, and the value of tags equals the exact gain from tailoring intensities across groups (Theorem~\ref{thm:partial-minimax}).
  \item \emph{Value of class-tagged disclosure and erosion under relabeling.} Class-tagged flow reports eliminate the spacing penalty for weight recovery (Theorem~\ref{thm:tags}), reducing the sample requirement from order $(\Delta\eta)^{-(4K-2)}$ to a spacing-free constant. However, when restricted traders can route activity through accounts carrying another label at convex resource cost $\gamma$, the optimal targeted restriction is partial, the equilibrium share of restricted activity that hides is independent of $\gamma$, and class tags lose all value as $\gamma \to 0$ (Proposition~\ref{prop:relabeling}). Only an owner-linked audit trail that ties accounts to beneficial owners retains its value.
  \item \emph{The surveillance paradox.} Without class tags, an untagged algorithm facing a level-$\alpha$ detection test crowds its memory rate toward its neighbor as the supervisor's data record $N$ grows. Whenever the identification divergence scales as $c N D^p$ with $p > 2$ (including $p=4$ for symmetric splitting and $p=6$ for moment-matched two-group hiding), crowding more than offsets the added data: the supervisor's equilibrium information falls to zero at rate $N^{-2/(p-2)}$ (e.g., $N^{-1/2}$ for $p=6$), detection probability collapses to the false-alarm rate $\alpha$, and the trader incurs a permanent strategy distortion cost (Theorem~\ref{thm:paradox}). Class tags eliminate this crowding.
\end{enumerate}

\textbf{What is new, and what is not.} The minimax rate $\Delta^{-(4K-2)}$ is the established rate for finite mixtures with near-colliding components (Heinrich and Kahn, 2018; Wu and Yang, 2020) and super-resolution of point sources (Batenkov et al., 2021), while the modulus characterization of linear functionals builds on Donoho (1994). We do not claim those statistical theorems as new. Our contribution is economic: we show that kernel sufficiency survives endogenous inventory-equilibrium pricing; establish the exact equivalence between supervisory identification and forecast loss in a feedback market; prove the minimax optimality of uniform (and group-uniform) supervisory policies; solve the strategic relabeling disclosure game comparing class tags to owner-linked audit trails; and derive the negative equilibrium identification rate of the surveillance paradox.

\textbf{Empirical illustration.} Using daily class-tagged institutional order flows from the Taiwan Stock Exchange (300 stocks, 2013--2016, 270,060 stock-days), we calibrate our theoretical bounds to observed institutional rules (foreign investors, investment trusts, and dealers). Tagged flows carry $4.7 \times 10^5$ times more information per stock-day about class composition than prices alone. Over all compositions statistically indistinguishable from the full panel price record, the effect of a uniform feedback cut is tightly pinned to $[0.0014, 0.0019]$, whereas the effect of halting foreign feedback spans $[-0.028, 0.028]$.

\subsection{Related literature}

\textbf{Aggregation and observational equivalence.} Exponential memory is classical in dynamic economics, spanning adaptive expectations (Cagan, 1956; Nerlove, 1958), geometric distributed lags (Koyck, 1954), and steady-state Kalman filtering (Muth, 1960). That aggregate dynamics depend only on a distribution of micro coefficients dates back to Granger (1980), Forni and Lippi (1997), Zaffaroni (2004), and Pesaran and Chudik (2014), with recovery of the cross-sectional distribution forming the disaggregation problem (Celov et al., 2010). Observational equivalence in policy evaluation traces to Sargent (1976) and, for sufficient statistics in heterogeneous-agent macroeconomics, Auclert et al.\ (2024). A large literature studies heterogeneous fundamentalist and chartist asset-pricing models (Brock and Hommes, 1998; Hommes, 2006) and estimates trader shares from prices (Boswijk et al., 2007; Franke and Westerhoff, 2012; Lux and Zwinkels, 2018); our results establish precisely which features of a trader population are recoverable from prices and at what statistical cost. Our decision-theoretic framing connects Blackwell (1953) and Le Cam (1964) with data-driven control (van Waarde et al., 2020) and aggregate forecasting (L\"utkepohl, 1987; Hendry and Hubrich, 2011).

\textbf{Microstructure, trader identity, and algorithmic supervision.} Prior microstructure studies show that broker or trader identity carries price-discovery information when order flow reflects private fundamentals (Frino et al., 2010; Linnainmaa and Saar, 2012), while Sanders et al.\ (2009) document weak predictive content of large-trader positions in agricultural futures. Our setting isolates a distinct mechanism: under a common public signal, identity is irrelevant for price discovery and forecasting, yet essential for assigning targeted supervisory interventions. In regulatory policy, Budish et al.\ (2015) advocate market-wide structural remedies (frequent batch auctions), Kirilenko et al.\ (2017) use account-level audit-trail data to reconstruct the Flash Crash, and Menkveld (2016) surveys high-frequency trading. SEC Rule~613 established the Consolidated Audit Trail (CAT) to link orders to customer identities across venues (Securities and Exchange Commission, 2012). Our disclosure game—building on costly state falsification (Lacker and Weinberg, 1989; Crocker and Morgan, 1998)—provides a formal rationale for that architecture: market-wide remedies require only prices, targeted restrictions require class records, and relabeling incentives necessitate owner-linked audit trails.

\textbf{Roadmap.} Section~\ref{sec:model} presents the baseline model and unified notation. Section~\ref{sec:sufficiency} establishes kernel sufficiency and noiseless identification. Section~\ref{sec:equilibrium} derives price impact and persistence in an inventory equilibrium and proves that all identification results carry over. Section~\ref{sec:cost} quantifies the sample complexity of composition versus market-wide questions. Section~\ref{sec:forecast} equates pairwise detection with forecast mismatch. Section~\ref{sec:supervision} analyzes minimax policy, disclosure, relabeling, the surveillance paradox, and the Taiwan calibration. Section~\ref{sec:conclusion} discusses scope and concludes. Appendix~\ref{app:proofs} collects self-contained proofs of the main economic theorems.

\section{Model and Notation}\label{sec:model}

Time is discrete, $t = 0, 1, 2, \dots$. Let $x_t$ denote the pricing error—the gap between the transaction price $p_t$ and the asset's fundamental value $F_t$. In the reduced-form benchmark, a dealer sector absorbs net order flow with price impact $\kappa > 0$, and arbitrage pulls the pricing error toward zero with persistence $a \in [0, 1)$. (Section~\ref{sec:equilibrium} derives both $a$ and $\kappa$ endogenously from competitive inventory-averse dealers.)

There are $K \ge 1$ groups of feedback traders with adaptation rates $\eta_i \in (0, 1)$ and activity weights $q_i$. Group $i$ maintains an exponentially weighted average $y_{i,t}$ of past pricing errors and submits feedback flow $q_i y_{i,t}$:
\begin{equation}\label{eq:baseline}
  x_t = a x_{t-1} + \kappa \sum_{i=1}^K q_i y_{i,t-1} + \kappa v_t + \varepsilon_t,
  \qquad
  y_{i,t} = (1 - \eta_i) y_{i,t-1} + \eta_i x_t,
\end{equation}
where $v_t$ is an observer's executed probe trade (with $v_t \equiv 0$ under passive observation) and $\varepsilon_t$ is a price innovation (news). As shown in the Online Supplement (Lemma~S1.1), both a trading desk tracking a latent random-walk trend via a steady-state Kalman filter (Muth, 1960) and a desk updating via $L_2$-regularized online gradient descent satisfy the exact recursion in~\eqref{eq:baseline}.

Write $r_i = 1 - \eta_i \in (0, 1)$ for group $i$'s \emph{memory rate} and $w_i = q_i \eta_i$ for its \emph{merged trading weight}. Define the \emph{merged measure} $M = \sum_{i=1}^K w_i \delta_{r_i}$, which sums weights over any groups sharing an identical memory rate. The \emph{feedback kernel} is the sequence of power moments of $M$:
\begin{equation}\label{eq:kernel}
  k_m = \sum_{i=1}^K w_i r_i^m = \int r^m \, dM(r), \qquad m = 0, 1, 2, \dots
\end{equation}
Total feedback flow is the causal convolution of past pricing errors with $k$. Define the \emph{feedback budget}
\begin{equation}\label{eq:budget}
  S = a + \kappa \sum_{i=1}^K q_i = a + \kappa \sum_{m=0}^\infty k_m.
\end{equation}
In the positive-feedback economy ($q_i \ge 0$), the closed-loop system~\eqref{eq:baseline} is Schur stable if and only if $S < 1$. Table~\ref{tab:notation} summarizes our unified notation across sections.

\begin{table}[!htbp]
\centering
\caption{Summary of core notation across sections.}
\label{tab:notation}
\small
\begin{tabularx}{\linewidth}{@{}lX@{}}
\toprule
Symbol & Economic meaning \\
\midrule
$x_t, F_t, p_t$ & Pricing error ($x_t = p_t - F_t$), public fundamental, and transaction price \\
$a, \kappa$ & Pricing-error persistence ($a \in [0,1)$) and price impact ($\kappa > 0$) \\
$\eta_i, r_i, q_i, w_i$ & Group $i$ adaptation rate, memory rate ($r_i = 1-\eta_i$), activity weight, merged weight ($w_i = q_i \eta_i$) \\
$M, k_m, G(z)$ & Merged measure $\sum_i w_i \delta_{r_i}$, feedback kernel $k_m = \int r^m dM$, Stieltjes transform $\sum_i w_i/(z - r_i)$ \\
$S, \bar{s}$ & Feedback budget $S = a + \kappa \sum_i q_i$ and prudential ceiling $\bar{s} \in (a, 1)$ \\
$H_t, \bar{a}, A$ & Dealer inventory, dealer carry share, and CARA risk aversion (Section~\ref{sec:equilibrium}) \\
$\Delta, \Delta\eta$ & Spread (half-width) of trader type parameters or memory rates around center $r_c$ \\
$u_d, U, \theta_c$ & Class $d$ feedback contribution ($u_d = \kappa q_d$), total feedback $U = \sum_d u_d$, policy scale $\theta_c$ \\
$\theta_{\mathcal{B}}, \theta_{\mathcal{H}}, \gamma, \alpha_{\mathcal{H}}$ & Retention rates for benign ($\mathcal{B}$) and harmful ($\mathcal{H}$) labels, relabeling cost $\gamma$, harm $\alpha_{\mathcal{H}} = -n_{\mathcal{H}}$ \\
$D, D_0, p, J$ & Distance to neighbor's memory, preferred distance, divergence exponent ($p > 2$), divergence $J = cND^p$ \\
\bottomrule
\end{tabularx}
\end{table}

\textbf{Observation regimes and the role of impact.} In the reduced-form identification analysis, the observer knows $(a, \kappa)$ and the noise law and observes $x_t$ (or a public signal from which $x_t$ is reconstructible). Giving the observer known $(a, \kappa)$ and exact observability of $x_t$ only strengthens the observer; hence all sample-complexity lower bounds hold \emph{a fortiori} when $(a, \kappa)$ must be estimated or when only raw transaction prices $p_t = F_t + x_t$ are observed with hidden fundamentals (Remark~\ref{rem:obs}). We analyze both passive price observation ($v_t = 0$) and active probing from a reset state. Reset probing is strictly an information-theoretic benchmark—not a proposed regulatory policy—demonstrating that even an unrealistically powerful probing supervisor cannot escape the $\Delta^{-(4K-2)}$ barrier. Finally, mechanical impact in~\eqref{eq:baseline} describes limit-order-book mechanics or inventory-averse dealers (Section~\ref{sec:equilibrium}); by contrast, a risk-neutral rational market maker who knows the aggregate flow rule nets out predictable feedback flow entirely, so prices reveal \emph{less}, not more, about the population (Niroula et al., 2026b, Proposition~C2.1).

\section{Kernel Sufficiency and Identification}\label{sec:sufficiency}

To establish sufficiency in full generality, embed the feedback population in a strategic market. An asset value $v$ is drawn, noise traders submit orders $z_t$, a strategic insider submits orders $X_t(v, \mathcal{H}_{t-1})$, and a dealer sets prices $P_t(\mathcal{H}_{t-1}, Y_t)$ from public history $\mathcal{H}_{t-1}$ and total order flow $Y_t = X_t + z_t + f_t$. Feedback groups update on a public signal $s_t = \mathcal{S}_t(p_0, \dots, p_t)$ and contribute aggregate flow $f_t = \sum_{m < t} k_m s_{t-1-m}$. The maps $X$, $P$, and $\mathcal{S}$ are arbitrary causal rules. The strategic players are the insider and the dealer; feedback-group rules are fixed, and their individual payoffs are not part of the payoff equivalence.

\begin{theorem}[The kernel is sufficient]\label{thm:sufficiency}
If two feedback-group configurations satisfy $k_m = k'_m$ for all $m < T$, then for every strategy profile $(X, P)$ and every realization of $(v, z)$, the price path, total order-flow path, insider trades and realized profit, and dealer inventory coincide on all dates $t \le T$. The two games share identical strategy sets and payoff functions, and hence identical equilibria under any solution concept defined from public histories and payoffs. The same holds for any finite population whose individual flows are causal linear functions of the public signal, with $k$ the sum of individual kernels. Group-level quantities, such as the weight of the largest group, are not invariant.
\end{theorem}

The proof is a short induction on dates (Appendix~\ref{app:proofs}). Its economic force lies in the object it isolates: a single population-level sequence $k$ settles every equilibrium concept simultaneously. Because the induction uses only that aggregate feedback flow is a fixed causal map of public signals, sufficiency extends to arbitrary nonlinear or profit-switching populations via the aggregate flow map (Online Supplement, Section~S2). Moreover, equal kernels remain observationally equivalent even when anonymous aggregate order flow is added to the public record.

The converse holds under the mechanical-impact observation experiment of Section~\ref{sec:model}:

\begin{theorem}[What prices identify]\label{thm:what-id}
In model~\eqref{eq:baseline}, fix known $(a, \kappa)$ and zero initial group states, and compare noiseless probe responses. Two economies produce identical price responses to every probe sequence if and only if their merged measures $M$ coincide.
\end{theorem}

Together, Theorems~\ref{thm:sufficiency} and~\ref{thm:what-id} establish that the feedback kernel (equivalently, the merged measure $M$, which has finitely many atoms) is a \emph{minimal identifying functional} of the trader population for prices: it determines everything prices can reveal, and noiseless prices determine it.

\section{Sufficiency and Identification in Inventory Equilibrium}\label{sec:equilibrium}

Section~\ref{sec:sufficiency} took price impact $\kappa$ and persistence $a$ as primitives. We now derive both endogenously from competitive, risk-averse dealers who carry inventory risk, proving that kernel sufficiency and the $(\Delta\eta)^{-(4K-2)}$ identification barrier hold in equilibrium.

\textbf{The inventory economy.} Consider an infinite-horizon discrete-time market:
\begin{itemize}
  \item \emph{Fundamental and price.} The asset's fundamental value $F_t$ follows a public random walk with i.i.d.\ innovations of variance $\sigma_F^2$. The transaction price is $p_t = F_t + x_t$, where $x_t$ is the equilibrium pricing error induced by dealer inventory absorption.
  \item \emph{Feedback groups.} Group $i$ submits orders $u_{i,t+1} = r_i u_{i,t} + w_i x_t$, with memory rates $r_i = 1 - \eta_i \in (0, 1)$ and merged weights $w_i = q_i \eta_i$.
  \item \emph{Exogenous flow.} Total customer order flow is $Y_t = \sum_{i=1}^K u_{i,t} + \xi_t$, where $\xi_t \sim N(0, \sigma_\xi^2)$ is i.i.d.\ noise trading (or supervisory probing priced by dealers as noise).
  \item \emph{Competitive dealers.} Overlapping generations of dealers with CARA risk aversion $A > 0$ absorb customer flow and lay off a fraction $1 - \bar{a} \in (0, 1]$ of inventory each period, so aggregate dealer inventory evolves as $H_t = \bar{a} H_{t-1} - Y_t$. Market clearing requires $H_t = (\E_t[p_{t+1}] - p_t)/(A V)$, where $V = \Var_t(p_{t+1})$ is the one-step conditional price variance.
\end{itemize}
Let $G(D) = \sum_{i=1}^K w_i / (D - r_i) = \int (D - r)^{-1} \, dM(r)$ denote the Stieltjes transform of the merged measure $M$.

\begin{theorem}[Price impact and persistence in inventory equilibrium]\label{thm:inv-eq}
Consider stationary linear equilibria in which $x_t$ is linear in dealer inventory $H_t$ and next-period feedback trades $\{u_{i,t+1}\}$.
\begin{enumerate}
  \item[(i)] \emph{Characterization.} The pricing law
  \[
    x_t = -\kappa H_t + \kappa \sum_{i=1}^K \frac{u_{i,t+1}}{D - r_i}, \qquad D \neq 0, \; D \neq r_i,
  \]
  is an equilibrium if and only if the induced state dynamics are Schur stable and $(\kappa, D, V)$ solve
  \begin{equation}\label{eq:eq-system}
    D = 1 - \kappa G(D), \qquad \kappa(D - \bar{a}) = A V D, \qquad V = \sigma_F^2 + \frac{\kappa^2 \sigma_\xi^2}{D^2}.
  \end{equation}
  An unexpected unit of order flow moves the price by $\kappa / D$.
  \item[(ii)] \emph{No-feedback benchmark.} When $w_i = 0$, $D = 1$ and $x_t = \bar{a} x_{t-1} + \kappa Y_t$. Persistence equals the dealer carry share $\bar{a}$, and impact solves $\kappa(1 - \bar{a}) = A(\sigma_F^2 + \kappa^2 \sigma_\xi^2)$. An equilibrium exists if and only if $\Lambda_0 \equiv (1 - \bar{a})^2 - 4 A^2 \sigma_F^2 \sigma_\xi^2 \ge 0$; when $\Lambda_0 > 0$, the low-impact root $\kappa_- \approx A \sigma_F^2 / (1 - \bar{a})$ is the equilibrium price of inventory risk.
  \item[(iii)] \emph{Invariance.} The equilibrium tuple $(\kappa, D, V)$ and the resulting price process depend on the trader population solely through $G$, and hence solely through the merged measure $M$.
  \item[(iv)] \emph{Small feedback.} For small aggregate feedback intensity $G(1) = \sum_i q_i$ with $\Lambda_0 > 0$, there exists a unique equilibrium near $(\kappa_-, 1)$ satisfying
  \[
    D = 1 - \kappa_- G(1) + o(G(1)), \qquad
    \kappa = \kappa_- + \frac{\kappa_-^2(\bar{a} + 2 A \kappa_- \sigma_\xi^2)}{\sqrt{\Lambda_0}}\, G(1) + o(G(1)).
  \]
  Positive feedback increases both equilibrium price impact $\kappa$ and unexpected shock impact $\kappa / D$.
\end{enumerate}
\end{theorem}

\begin{remark}[Observability of pricing errors]\label{rem:obs}
Dealers and feedback traders in Theorem~\ref{thm:inv-eq} observe $x_t = p_t - F_t$ because $F_t$ is public among market participants while inventory-induced mispricing $x_t$ is earned by risk-averse dealers. An outside supervisor who observes only raw transaction prices $p_t = F_t + x_t$ without observing the random-walk fundamental $F_t$ receives strictly less information. Every sample-complexity lower bound derived below therefore holds \emph{a fortiori} under unobserved fundamentals (while Corollary~S7.2 in the Online Supplement shows that unknown primitives do not alter the upper-bound rate under active probes).
\end{remark}

\begin{remark}[Mispricing feedback versus raw-return feedback]\label{rem:returns}
Theorem~\ref{thm:inv-eq} specifies feedback rules on the pricing error $x_t = p_t - F_t$, which isolates pure inventory-driven feedback in closed form. If algorithmic traders instead condition on raw price changes $\Delta p_t = \Delta F_t + \Delta x_t$, fundamental innovations $\Delta F_t$ enter the traders' memory states directly alongside inventory shocks. This introduces a second exogenous innovation channel (analyzed in Online Supplement Section~S8) without altering pathwise kernel sufficiency (Theorem~\ref{thm:sufficiency}) or relaxing the $\Delta^{-(4K-2)}$ identification barrier for closely spaced memory rates.
\end{remark}

\begin{theorem}[Identification results carry over to equilibrium]\label{thm:eq-id}
In any stationary equilibrium of Theorem~\ref{thm:inv-eq} with $D \neq \bar{a}$:
\begin{enumerate}
  \item[(i)] \emph{Exact reduced form.} Path by path, the equilibrium pricing error follows the closed loop~\eqref{eq:baseline} with effective persistence $a^* = \bar{a}/D$, effective impact $\kappa^* = \kappa/D$, unchanged memory rates $r_i$, and rescaled merged weights $\tilde{w}_i = w_i (D - \bar{a})/(D - r_i)$.
  \item[(ii)] \emph{Transmission zeros.} The memory rates $r_i$ remain the exact transmission zeros of the equilibrium price response.
  \item[(iii)] \emph{Identification boundary.} The effective merged measure $\tilde{M} = h M$, where $h(r) = (D - \bar{a})/(D - r)$ is positive whenever $D > \max_i r_i$, is in one-to-one correspondence with the primitive merged measure $M$. Thus $M$ remains the exact identifying functional of equilibrium prices.
  \item[(iv)] \emph{Minimax rate preservation.} Under the uniform stability condition $a^* + \kappa^* W / (1 - r_c - \Delta\eta) < 1$ of Theorem~\ref{thm:cost}, two inventory equilibria with identical primitives $(A, \sigma_F^2, \sigma_\xi^2, \bar{a})$ that differ only in their feedback-trader composition require $N \ge c \sigma^2 (\Delta\eta)^{-(4K-2)}$ observations to be distinguished.
\end{enumerate}
\end{theorem}

The proofs of Theorems~\ref{thm:inv-eq} and~\ref{thm:eq-id} are given in Appendix~\ref{app:proofs}. Economically, forward-looking dealers anticipate predictable feedback trades and incorporate them into current quotes; rather than obscuring the memory rates, anticipation leaves $r_i$ untouched and up-weights slower-memory groups by $(D - \bar{a})/(D - r_i)$.

To organize the decision-theoretic consequences of sufficiency, we state the following immediate corollary of Blackwell (1953) and Le Cam (1964):

\begin{proposition}[What an intervention requires]\label{prop:decision-suff}
Suppose supervisory observations depend on the trader population only through its aggregate flow map.
\begin{enumerate}
  \item[(i)] If the loss of an action depends on the population only through the flow map, the set of safe actions is identified by prices.
  \item[(ii)] If two populations share the same flow map and an action is safe in one but unsafe in the other, every observation-based rule takes that action with identical probability in both.
  \item[(iii)] If the observation laws of two populations lie within total-variation distance $\delta$, any rule that takes the action with probability $p$ where it is safe takes it with probability at least $p - \delta$ where it is unsafe.
\end{enumerate}
\end{proposition}

\section{The Sample-Complexity Price of Composition}\label{sec:cost}

Theorem~\ref{thm:what-id} establishes that groups with distinct memory rates are identifiable in principle. We now quantify how much data identification actually requires when memory rates are closely spaced, comparing passive price observation against the active-probing benchmark.

\begin{theorem}[The cost of composition]\label{thm:cost}
Fix $K \ge 2$, a center memory rate $r_c \in (0,1)$, total merged weight $W > 0$, and spacing $\Delta\eta > 0$ such that every candidate economy is Schur stable.
\begin{enumerate}
  \item[(i)] \emph{Passive price observation.} When prices are observed passively as the stationary output of~\eqref{eq:baseline} driven by i.i.d.\ Gaussian news $\varepsilon_t \sim N(0, \sigma^2)$, there exist two $K$-group economies with rates in $r_c + \Delta\eta[-1,1]$ and positive weights whose merged measures match all moments of order $0, \dots, 2K-2$, and whose Kullback--Leibler divergence per observation is of exact order $(\Delta\eta)^{4K-2}$. Distinguishing them with both error probabilities below $1/4$ requires $T \ge c (\Delta\eta)^{-(4K-2)}$ passive observations.
  \item[(ii)] \emph{Active probing is no cheaper.} Even if the supervisor can run $N$ independent reset experiments recording the noisy impulse response with Gaussian observation noise of variance $\sigma^2$, the same pair has impulse-response gap of order $(\Delta\eta)^{2K-1}$, requiring $N \ge c \sigma^2 (\Delta\eta)^{-(4K-2)}$ experiments.
  \item[(iii)] \emph{Matching upper bound.} On any compact class of $\gamma$-separated rescaled rates with weights bounded away from zero, Prony's method applied to the first $2K$ estimated kernel values recovers the weights from $N = C \sigma^2 (\Delta\eta)^{-(4K-2)}$ reset experiments.
\end{enumerate}
\end{theorem}

The proofs of parts (i)--(iii) appear in the Online Supplement (Theorems~S4.1, S4.2, and~S4.4). The hard pair consists of the $K$-point Gauss--Legendre and Gauss--Radau quadrature rules scaled to $r_c + \Delta\eta[-1,1]$, which match all moments through degree $2K-2$ while differing at order $2K-1$. For two groups ($K=2$), halving the memory spread multiplies the required data by $2^6 = 64$; for three groups ($K=3$), by $2^{10} = 1{,}024$. Gaussian news is the least-favorable innovation law: for any regular non-Gaussian news law whose location divergence is quadratically bounded (such as Laplace innovations), the exponent $4K-2$ is unchanged (Online Supplement, Section~S4).

\textbf{The barrier does not lift near instability.} One might conjecture that as a market approaches instability ($S \uparrow 1$), feedback amplification magnifies compositional differences and makes traders easier to distinguish. The opposite holds: at a common feedback budget $S = 1 - \epsilon$, the passive KL divergence per observation between the hard pair of Theorem~\ref{thm:cost} remains bounded by $C(\Delta\eta)^{4K-2}$ uniformly in $\epsilon > 0$ (Online Supplement, Proposition~S4.5). As Figure~\ref{fig:near-instability} illustrates, a market near instability reveals how close it is to instability, not who is trading in it. Table~\ref{tab:years} translates these lower bounds into illustrative calendar years of daily price data.

\begin{figure}[!htbp]
  \centering
  \includegraphics[width=\linewidth]{fig1.png}
  \caption{The passive identification barrier near instability. Both panels plot the Gaussian Kullback--Leibler divergence per observation between the stationary price processes of the $K$-node Gauss and Gauss--Radau populations spread over $\Delta\eta$ around $r_c = 0.5$ (total weight $0.3$, impact $\kappa = 0.5$). Panel~(a): at a common feedback budget $1 - S = 10^{-6}$, information decays as $(\Delta\eta)^{4K-2}$ (slopes $6$ and $10$). Panel~(b): holding the feedback budget common across economies (solid curves), information stays uniformly bounded as $1 - S \to 0$; holding persistence $a$ common instead (dashed curves), information diverges only because the two economies' feedback budgets then differ.}
  \label{fig:near-instability}
\end{figure}

\begin{table}[!htbp]
\centering
\caption{Illustrative years of daily price data below which no test separates the two populations of Theorem~\ref{thm:cost}(i) with both error probabilities at $5\%$.}
\label{tab:years}
\small
\begin{tabular}{@{}llccc@{}}
\toprule
Groups & Memory rates & Half-lives (days) & $1 - S = 0.1$ & $1 - S = 10^{-6}$ \\
\midrule
$K = 2$ & $0.5 \pm 0.1$   & $0.8\text{--}1.4$  & $20{,}000$        & $16{,}000$ \\
        & $0.5 \pm 0.05$  & $0.9\text{--}1.2$  & $1.2 \times 10^6$ & $9.0 \times 10^5$ \\
        & $0.8 \pm 0.1$   & $1.9\text{--}6.6$  & $100$             & $69$ \\
        & $0.8 \pm 0.05$  & $2.4\text{--}4.3$  & $6{,}000$         & $3{,}900$ \\
        & $0.9 \pm 0.05$  & $4.3\text{--}13.5$ & $120$             & $77$ \\
        & $0.9 \pm 0.025$ & $5.2\text{--}8.9$  & $7{,}000$         & $4{,}300$ \\
\midrule
$K = 3$ & $0.5 \pm 0.1$   & $0.8\text{--}1.4$  & $1.5 \times 10^8$ & $9.9 \times 10^7$ \\
        & $0.5 \pm 0.05$  & $0.9\text{--}1.2$  & $1.5 \times 10^{11}$ & $9.6 \times 10^{10}$ \\
\bottomrule
\end{tabular}
\vspace{4pt}
\parbox{\linewidth}{\footnotesize\emph{Notes:} Lower bound $T \ge \ln(5)/\KL$ (Bretagnolle--Huber), computed from the closed-form Gaussian spectra at common feedback budget $S$, $\kappa = 0.5$, and feedback share $\kappa \sum_i q_i = 0.3$, assuming $252$ trading days per year. For 5-minute intraday memory horizons recorded continuously, divide by $417$. Entries are worst-case lower bounds for moment-matched pairs, not fitted venue horizons.}
\end{table}

\subsection{Which questions prices answer cheaply}

Composition is only one functional of the merged measure $M$. Consider any linear functional $F_g(M) = \int g(r)\, dM(r)$ for continuous $g$ on $[0, \bar{r}]$ ($\bar{r} < 1$), with total mass at most $\bar{W}$, observed through the first $L$ kernel taps with $N(0, \sigma^2)$ noise across $N$ experiments. Write $\mathcal{R}(h) = (\sup h)_+ + (-\inf h)_+$ for the signed oscillation of $h$.

\begin{theorem}[Which questions prices answer cheaply]\label{thm:cheap}
Let $z_\beta$ denote the standard normal $\beta$-quantile and $\alpha \in (0, 1/2)$.
\begin{enumerate}
  \item[(i)] \emph{Exact modulus cost.} The minimum number of experiments $N_{\mathrm{cert}}(\epsilon; F_g)$ required to certify $F_g(M)$ to half-width $\epsilon$ at confidence level $1-\alpha$ is pinned, up to the quantile ratio $\lambda_\alpha = z_{1-\alpha/2}/z_{1-\alpha}$ (approximately $1.19$ at the $95\%$ level), by the polynomial approximation problem
  \[
    \omega_L(\delta) = \min_{c \in \R^L} \Bigl\{ \delta \|c\|_2 + \bar{W}\, \mathcal{R}\Bigl(g - \sum_{t < L} c_t r^t\Bigr) \Bigr\}
  \]
  evaluated at $\delta = 2 z_{1-\alpha} \sigma / \sqrt{N}$.
  \item[(ii)] \emph{Cheap (parametric) questions.} $N_{\mathrm{cert}}(\epsilon; F_g) = O(\sigma^2 \epsilon^{-2})$—independent of the number of groups $K$ and their memory spacing—if and only if $g$ agrees on $[0, \bar{r}]$ with a power series having square-summable coefficients.
  \item[(iii)] \emph{Impossible questions.} Without a memory bound $\bar{r} < 1$, any functional satisfying
  \[
    \limsup_{r \uparrow 1}\, |g(r)|\sqrt{1 - r} = \infty
  \]
  cannot be certified to any finite accuracy from any finite number of experiments.
\end{enumerate}
\end{theorem}

Theorem~\ref{thm:cheap} (Online Supplement, Theorem~S10.1 and Proposition~S10.2) specializes Donoho's (1994) modulus of continuity to our feedback market. It cleanly separates market-wide questions from compositional ones: a single kernel tap and discounted feedback $\int (1 - \beta r)^{-1} dM$ ($\beta < 1$) are parametric ($O(\sigma^2 \epsilon^{-2})$); total feedback weight $\Phi = \sum_m k_m$ sits on the boundary at $O(\sigma^2 \epsilon^{-2} \log(1/\epsilon))$ (Niroula et al., 2026a); and individual group weights are nonlinear functionals costing $(\Delta\eta)^{-(4K-2)}$.

\begin{corollary}[Separation of market-wide and compositional inference]\label{cor:sep}
Fix $K \ge 2$, accuracy $\epsilon > 0$, and any cheap functional $F_g$ satisfying Theorem~\ref{thm:cheap}(ii). Over $K$-group populations with memory spacing $\Delta\eta$, the ratio of experiments needed to recover group weights within a fixed fraction of total weight to the experiments needed to certify $F_g$ within $\epsilon$ is at least $c \epsilon^2 (\Delta\eta)^{-(4K-2)}$, diverging as $\Delta\eta \to 0$.
\end{corollary}

\subsection{Beyond exponential memory: smooth rule families}

The exponent $4K-2$ does not rely on exponential memory or linear trading rules. Index trader types by $\theta \in \Theta \subset \R$ with arbitrary causal rules $f_\theta(h)$ mapping public history $h = (x_t, x_{t-1}, \dots)$ to order flow, such as power-law memory, hump-shaped lags, fixed lookback windows, or saturating nonlinear rules. A population is a finite measure $\nu$ on $\Theta$ generating aggregate flow $F_\nu(h) = \int f_\theta(h)\, \nu(d\theta)$, entering prices via~\eqref{eq:baseline} with $\kappa F_\nu(h_{t-1})$ in place of the linear feedback term.

\begin{theorem}[Lower bound for smooth families]\label{thm:smooth}
Fix $K \ge 2$, center $\theta_c$, total weight $W > 0$, and spread $\Delta > 0$ with $[\theta_c - \Delta, \theta_c + \Delta] \subset \Theta$. Suppose $\theta \mapsto f_\theta(h)$ has $2K-1$ continuous derivatives at every history with $|\partial_\theta^{2K-1} f_\theta(h)| \le B(h)$, and conditional on past observations and probes, price innovations are independent $N(0, \sigma^2)$. Compare records from a common observed initial history under any predictable probe policy.
\begin{enumerate}
  \item[(i)] \emph{Flow gap.} There exist two $K$-type populations $\nu, \nu'$ supported on $[\theta_c - \Delta, \theta_c + \Delta]$ with positive weights summing to $W$ and Wasserstein distance $W_1(\nu, \nu') = d_K W \Delta$ ($d_K > 0$) such that, at every history $h$,
  \[
    |F_\nu(h) - F_{\nu'}(h)| \le \frac{2W}{(2K-1)!}\, B(h)\, \Delta^{2K-1}.
  \]
  \item[(ii)] \emph{Lower bound.} If $\E_\nu[B(h_t)^2] \le \bar{B}^2 < \infty$ uniformly across dates (and, for asymptotic claims, uniformly over spreads and policies), distinguishing $\nu$ from $\nu'$ with both error probabilities below $1/4$ requires at least
  \[
    T \ge \frac{((2K-1)!)^2}{4 \kappa^2 W^2 \bar{B}^2}\, \sigma^2 \Delta^{-(4K-2)}
  \]
  periods.
\end{enumerate}
\end{theorem}

The self-contained proof in Appendix~\ref{app:proofs} applies Taylor's remainder theorem to the Gauss--Legendre and Gauss--Radau quadrature pair and invokes the sequential Gaussian KL chain rule, establishing the paper's headline $\Delta^{-(4K-2)}$ sample-complexity lower bound directly. For general linear memory kernels satisfying a local derivative rank condition, a matching upper bound holds within a shrinking cluster:

\begin{theorem}[Recovery within a shrinking cluster of linear rules]\label{thm:recovery}
Fix known $K \ge 2$, $a$, $\kappa > 0$, $\theta_c$, $L \ge 2K$, and the compact separated class
\[
  \mathcal{Z} = \{(s, w) : -1 \le s_1 < \dots < s_K \le 1,\; s_{i+1} - s_i \ge \gamma,\; \omega_{\min} \le w_i \le \omega_{\max}\},
\]
with $\gamma > 0$ and $0 < \omega_{\min} \le \omega_{\max} < \infty$. Let $g = (g_0, \dots, g_{L-1})$ be $C^{2K}$ near $\theta_c$ with derivative vectors of orders $0, \dots, 2K-1$ linearly independent at $\theta_c$, and let $k_m = \sum_i w_i g_m(\theta_c + \Delta s_i)$ with known $\Delta > 0$. From $N$ independent copies of the deterministic unit-impulse response $(h_0, \dots, h_L)$ plus independent $N(0, \sigma^2 I_{L+1})$ noise, there exist $\Delta_0 > 0$ and $C < \infty$ such that for all $0 < \Delta \le \Delta_0$, every exact minimum-distance estimator $\hat{z}$ over $\mathcal{Z}$ satisfies, with probability at least $1 - \delta$,
\[
  \|\hat{z} - z\| \le C \sigma \bigl(\sqrt{L+1} + \sqrt{2\log(1/\delta)}\bigr)\, N^{-1/2} \Delta^{-(2K-1)}.
\]
\end{theorem}

The rank condition holds for both exponential memory $g_m(\theta) = (1-\theta)\theta^m$ ($\theta_c \in (0,1)$) and power-law memory $g_m(\theta) = (m+1)^{-\theta}$ ($\theta_c > 1$), matching the $\Delta^{-(4K-2)}$ sample order of Theorem~\ref{thm:smooth} (Online Supplement, Section~S18).

\section{What a Record Reveals: Identification Equals Forecast Loss}\label{sec:forecast}

Why can a supervisor calibrate a uniform feedback cut from prices while remaining unable to attribute feedback to individual classes? The next theorem provides the exact information-theoretic link: in a Gaussian feedback market, the information a price record carries about two competing populations is literally proportional to the cumulative forecast loss from mistaking one for the other.

\begin{theorem}[Pairwise detection and forecast mismatch]\label{thm:forecast-kl}
Consider two causal aggregate flow maps $F, F'$ under common known $(a, \kappa)$, common observed initial information, independent $N(0, \sigma^2)$ innovations, and a common predictable probe policy. For a record of $T$ subsequent periods, define the one-step excess mean-squared forecast error and cumulative signal-to-noise ratio
\[
  V_t = \kappa^2 \E_P\bigl[(F(h_t) - F'(h_t))^2\bigr], \qquad \Lambda_T = \sigma^{-2} \sum_{t < T} V_t.
\]
Then $V_t$ is the exact excess mean-squared error of forecasting $x_{t+1}$ using the alternative conditional mean under the true law $P$, and
\[
  \KL(P_T \| P'_T) = \frac{\Lambda_T}{2}.
\]
Both testing errors can lie below $1/4$ only if $\Lambda_T > 1$. Moreover, for any countable family of candidate densities with prior mass $\pi_* > 0$ on the truth, the Bayesian mixture $q$ satisfies the pathwise regret bound $\log p_*(X_{1:T}) - \log q(X_{1:T}) \le \log(1/\pi_*)$.
\end{theorem}

The proof is given in Appendix~\ref{app:proofs} (combining the Gaussian chain rule with the mixture likelihood bound of Dawid, 1984 and Vovk, 2001). Furthermore, composition is unnecessary for forecasting the pricing error $x_t$ itself: with known parameters, the price record and the class-flow record yield the exact same optimal forecast, and for an unknown stable population, a fitted autoregression of order $L_T = \lceil c \log T \rceil$ attains the oracle risk up to $O(\sigma^2 \log T / T)$ (Ing and Wei, 2005), while a Bayesian $\mathrm{AR}(L)$ achieves uniform cumulative log-loss $O(\log^2(T+1)/T)$ across all stable populations (Online Supplement, Theorems~S18.1 and~S18.2). What prices withhold from a supervisor is therefore not forecast accuracy, but the class attribution required for targeted policy.

\section{Supervision, Minimax Policy, and Strategic Disclosure}\label{sec:supervision}

We now derive the regulatory implications of observational equivalence, moving from static prudential ceilings to minimax policy design, class-tagged disclosure, strategic relabeling, endogenous crowding, and an empirical calibration.

\subsection{Market-wide versus targeted interventions}

Instantiating Proposition~\ref{prop:decision-suff}, suppose a supervisor must keep the feedback budget $S = a + \kappa \sum_i q_i$ below a prudential ceiling $\bar{s} \in (a, 1)$ in a positive, stable market with known $(a, \kappa)$.

\begin{proposition}[Market-wide versus targeted action]\label{prop:interventions}
\begin{enumerate}
  \item[(i)] \emph{Market-wide scaling.} Scaling all feedback activity uniformly by $\theta \in [0, 1]$ meets the ceiling if and only if $a + \theta(S - a) \le \bar{s}$, which depends on the population solely through the aggregate budget $S$ (and thus the kernel).
  \item[(ii)] \emph{Targeted halt.} Halting the largest group does not. For every $K \ge 2$, there exist two $K$-group economies whose merged measures match moments $0, \dots, 2K-2$ (so their aggregate budgets differ only by $O((\Delta\eta)^{2K-1})$) while their largest groups differ by a fixed positive margin. Over a nonempty interval of ceilings $\bar{s}$, halting the largest group satisfies the ceiling in one economy and breaches it in the other; deciding between them from $N$ reset records with both errors at most $1/4$ requires $N \ge c \sigma^2 (\Delta\eta)^{-(4K-2)}$.
\end{enumerate}
\end{proposition}

\begin{corollary}[Expected loss from weakly identified targeting]\label{cor:loss}
In the hard pair of Proposition~\ref{prop:interventions}, let economy~$0$ be the case where halting the largest class breaches $\bar{s}$ (incurring loss $b > 0$) while a uniform cut meets it, and economy~$1$ the case where both meet $\bar{s}$ (with unnecessary uniform restriction incurring loss $c > 0$). Under equal priors and observation laws $P_0, P_1$, every decision rule has average loss
\[
  R \ge \frac{\min(b, c)}{2}\bigl\{1 - \TV(P_0, P_1)\bigr\} \ge \frac{\min(b, c)}{2}\Bigl[1 - \sqrt{\KL(P_0 \| P_1)/2}\Bigr]_+.
\]
Under prices with $\KL(P_0 \| P_1) \le C T \Delta^{4K-2}$, achieving $R \le \epsilon \min(b, c)$ ($\epsilon < 1/2$) requires $T \ge 2(1 - 2\epsilon)^2 / (C \Delta^{4K-2})$. Under $N$ independent Gaussian class-tagged flow measurements with covariance $\tau^2 I$ and mean separation $d \neq 0$, the likelihood-ratio rule achieves $R_{\mathrm{tag}} \le \frac{1}{2}\max(b, c)\exp\{-N \|d\|^2 / (8\tau^2)\}$.
\end{corollary}

Figure~\ref{fig:intervention} illustrates this contrast in a two-class economy: halting the largest class produces qualitatively opposite prudential outcomes across nearly indistinguishable economies, whereas a single verified class report collapses decision error by six orders of magnitude in sample size.

\begin{figure}[!htbp]
  \centering
  \includegraphics[width=\linewidth]{fig2.png}
  \caption{Intervention choice and the information value of class reports. Panel~A uses the two-class construction of Online Supplement Section~S17 ($a = 0.3$, $\kappa = 0.5$, $W = 0.2$, $r_c = 0.5$): a ceiling of $0.375$ separates the post-halt budgets while uniform scaling satisfies it in both economies. Panel~B fixes $\Delta = 0.05$ and compares exact likelihood-ratio error probabilities across $N$ independent repetitions (30 impulse-response coordinates vs.\ one verified class report, both with Gaussian measurement standard deviation $0.01$): achieving $5\%$ error requires roughly $2.03 \times 10^6$ public-response repetitions versus $2$ class reports.}
  \label{fig:intervention}
\end{figure}

\subsection{Minimax optimality of uniform policies}

Let class $d \in \{1, \dots, K\}$ contribute $u_d = \kappa q_d$ to the feedback budget, with total feedback $U = \sum_{d=1}^K u_d$. A supervisory policy scales the class carrying label $c$ by $\theta_c \in [0, 1]$ at implementation cost $C(\theta)$. If $\sigma$ is the unknown permutation assigning labels to underlying classes, the post-policy budget is $S'_\sigma = a + \sum_{c=1}^K \theta_c u_{\sigma(c)}$ and the prudential loss is $\ell(S'_\sigma)$.

\begin{theorem}[Without tags, uniform policies are minimax]\label{thm:minimax}
Let $\ell$ be convex and $C$ be convex and permutation-invariant on $[0, 1]^K$.
\begin{enumerate}
  \item[(i)] \emph{Weak dominance.} For every policy $\theta \in [0, 1]^K$, the uniform policy $\bar{\theta}\mathbf{1}$ with $\bar{\theta} = K^{-1}\sum_{c=1}^K \theta_c$ yields budget $a + \bar{\theta} U$ independent of $\sigma$ and satisfies
  \[
    \ell(a + \bar{\theta} U) + C(\bar{\theta}\mathbf{1})
    \le \frac{1}{K!} \sum_{\sigma} \ell(S'_\sigma) + C(\theta)
    \le \max_{\sigma} \ell(S'_\sigma) + C(\theta).
  \]
  Because no price record distinguishes labelings (Theorem~\ref{thm:sufficiency}), every price-based policy is weakly dominated by a uniform policy against the worst consistent labeling.
  \item[(ii)] \emph{Uniqueness.} If, in addition, $C$ is strictly convex on $[0, 1]^K$ (or $\ell$ is strictly convex on $[a, a + U]$ with $U > 0$ and not all $u_d$ equal), then $\bar{\theta}\mathbf{1}$ strictly dominates every non-uniform policy $\theta \neq \bar{\theta}\mathbf{1}$, and the minimax problem admits a unique solution $\bar{\theta}^*\mathbf{1}$: the uniform cut is the unique minimax supervisory policy.
\end{enumerate}
\end{theorem}

\begin{corollary}[Guarantee of a single halt]\label{cor:halt}
In the setting of Theorem~\ref{thm:minimax}:
\begin{enumerate}
  \item[(i)] Halting any single label $F$ is guaranteed to reduce the feedback budget only by $\min_d u_d$, and some labeling consistent with the price record attains this worst case.
  \item[(ii)] A uniform cut to $\theta = (\bar{s} - a)/(S - a)$ achieves the ceiling $\bar{s}$ exactly in every population with budget $S$, and no policy meeting the ceiling removes less aggregate feedback than $S - \bar{s}$.
\end{enumerate}
\end{corollary}

Verified partial tags break permutation symmetry across tagged groups while preserving it within each group:

\begin{theorem}[Minimax value of partial tags]\label{thm:partial-minimax}
Suppose verified tags partition the classes into groups $g = 1, \dots, G$, so the record identifies each class's group membership but cannot distinguish relabelings within any group. If $\ell$ is convex and $C$ is convex and invariant to within-group permutations, then replacing $\theta$ by the group-average policy $\bar{\theta}_g = |g|^{-1} \sum_{c \in g} \theta_c$ yields budget $a + \sum_{g=1}^G \bar{\theta}_g U_g$ ($U_g = \sum_{d \in g} u_d$) under every consistent labeling and weakly dominates $\theta$ in minimax loss. The minimax value of tags is
\[
  V_{\mathrm{tag}} = \min_{\bar{\theta}} \max_{\sigma} \{\ell(S'_\sigma) + C\} - \min_{(\bar{\theta}_g)_g} \max_{\sigma} \{\ell(S'_\sigma) + C\} \ge 0,
\]
the exact gain from tailoring policy intensities across verified groups. As in Theorem~\ref{thm:minimax}(ii), if $C$ is strictly convex on $[0, 1]^K$, the group-uniform minimizer is unique.
\end{theorem}

The proofs of Theorems~\ref{thm:minimax} and~\ref{thm:partial-minimax} and Corollary~\ref{cor:halt} are short and self-contained (Appendix~\ref{app:proofs}), and their finite combinatorial statements are verified in Lean~4 (the Online Supplement documents exact formalization scope across all results, separating machine-checked deterministic steps from paper-level probabilistic arguments). In numerical checks across $14{,}400$ random economies with three convex losses and three cost functions, the uniform policy is strictly better than non-uniform policies in $86\%$ of cases; across $1{,}200$ random tag partitions, partial tags deliver strictly positive value in $27\%$ of economies.

\subsection{Class-tagged disclosure, relabeling, and the surveillance paradox}

\begin{theorem}[Statistical value of class-tagged disclosure]\label{thm:tags}
Suppose that in addition to prices, the supervisor receives each class's order flow with independent Gaussian measurement noise. On the rich class of Theorem~\ref{thm:cost}(iii), the number of experiments required to recover group weights drops from order $(\Delta\eta)^{-(4K-2)}$ to a constant independent of $\Delta\eta$, and for memory rates it drops to order $(\Delta\eta)^{-2}$.
\end{theorem}

As shown in the Online Supplement (Theorem~S5.4), tagging $m$ of $K$ groups lowers the composition exponent by $4$ per tag (to $4(K-m)-2$), so tagging all but one class removes the spacing penalty for weights entirely.

Crucially, however, disclosure regimes operate on strategic traders. Before solving the strategic relabeling and surveillance-crowding games, the next lemma anchors their reduced-form payoffs directly in the inventory equilibrium of Section~\ref{sec:equilibrium} and the steady-state Kalman trading desk of Section~\ref{sec:model}.

\begin{lemma}[Equilibrium microfoundation of disclosure and crowding payoffs]\label{lem:micro}
\begin{enumerate}
  \item[(i)] \emph{First-order welfare from inventory-equilibrium price impact.} In the small-feedback equilibrium of Theorem~\ref{thm:inv-eq}(iv), each unit of aggregate feedback intensity $G(1) = \sum_j q_j$ raises equilibrium impact $\kappa$ by $\kappa_-^2(\bar{a} + 2A\kappa_-\sigma_\xi^2)/\sqrt{\Lambda_0} > 0$, creating a marginal dealer-inventory and price-dislocation externality $\lambda_{\mathrm{ext}} > 0$. If class $j \in \{\mathcal{H}, \mathcal{B}\}$ earns private value $b_j > 0$ per retained unit and imposes marginal feedback externality $\lambda_{\mathrm{ext}} \phi_j$, then the first-order net social value of a retained unit is $n_j = b_j - \lambda_{\mathrm{ext}} \phi_j$, yielding $n_{\mathcal{H}} < 0 < n_{\mathcal{B}}$ whenever $b_{\mathcal{H}} < \lambda_{\mathrm{ext}}\phi_{\mathcal{H}}$ and $b_{\mathcal{B}} > \lambda_{\mathrm{ext}}\phi_{\mathcal{B}}$.
  \item[(ii)] \emph{Quadratic strategy-distortion cost from Kalman filtering.} Consider a trading desk tracking a latent random-walk trend $v_{t+1} = v_t + \zeta_{t+1}$ ($\Var(\zeta) = \sigma_\zeta^2 > 0$) from noisy signals $v_t + \nu_t$ ($\Var(\nu) = \sigma_\nu^2 > 0$) with exponential memory rate $r \in (0, 1)$, submitting the one-period optimal trade $u_t = \hat{v}_t / (\kappa + \lambda)$ that maximizes $\E[v_t u - \frac{1}{2}(\kappa + \lambda)u^2 \mid \hat{v}_t]$. Its stationary forecast-error variance is
  \[
    P(r) = \frac{\sigma_\zeta^2 + (1 - r)^2 \sigma_\nu^2}{1 - r^2}
         = \frac{\sigma_\zeta^2}{1 - r^2} + \frac{2\sigma_\nu^2}{1 + r} - \sigma_\nu^2,
  \]
  which is strictly convex on $(0, 1)$ with a unique minimizer $r_0 \in (0, 1)$ and $P''(r_0) > 0$. Deviating from $r_0$ to $r$ incurs per-period expected profit loss
  \[
    \Pi(r_0) - \Pi(r) = \frac{P(r) - P(r_0)}{2(\kappa + \lambda)}
    = k (r - r_0)^2 + O\!\left(|r - r_0|^3\right),
    \qquad k = \frac{P''(r_0)}{4(\kappa + \lambda)} > 0.
  \]
  For a small trader group ($w_i \ll 1$), shifting $r_i$ alters the market-wide Stieltjes transform $G(D)$—and hence equilibrium impact $\kappa$ in Theorem~\ref{thm:inv-eq}(iv)—only at order $O(w_i)$, so the quadratic distortion cost $k(D_0 - D)^2$ is exact at leading order.
\end{enumerate}
\end{lemma}

Class tags attached to accounts prevent a trader from hiding merely by tweaking its algorithm, but a restricted trader can still route orders through accounts carrying an unrestricted label.

\begin{proposition}[Disclosure when labels can be gamed]\label{prop:relabeling}
Fix one supervisory cycle. A share $\chi \in (0, 1)$ of feedback activity belongs to a harmful class $\mathcal{H}$ and $1 - \chi$ to a benign class $\mathcal{B}$. Each retained unit has net social value $n_{\mathcal{H}} < 0 < n_{\mathcal{B}}$, and $\mathcal{H}$ traders earn private gain $b_{\mathcal{H}} > 0$ per retained unit. Without disclosure, the supervisor chooses a uniform retention $\theta \in [0, 1]$. With class tags, the supervisor sets retention $\theta_{\mathcal{H}}$ for label $\mathcal{H}$ and $\theta_{\mathcal{B}} \ge \theta_{\mathcal{H}}$ for label $\mathcal{B}$, and $\mathcal{H}$ traders may route a share $z \in [0, 1]$ of their activity through $\mathcal{B}$ accounts at resource cost $\gamma z^2 / 2$ per unit ($\gamma > 0$). An owner-linked audit trail prevents relabeling ($z \equiv 0$). Write $\alpha_{\mathcal{H}} = -n_{\mathcal{H}} > 0$, $m = \chi n_{\mathcal{H}} + (1 - \chi) n_{\mathcal{B}}$, $s_{\mathcal{B}} = (1 - \chi) n_{\mathcal{B}}$, and $g = \gamma / \{b_{\mathcal{H}}(2\alpha_{\mathcal{H}} + b_{\mathcal{H}})\}$.
\begin{enumerate}
  \item[(i)] \emph{Optimal hiding.} $\mathcal{H}$ traders route $z^* = \min\{1,\, b_{\mathcal{H}}(\theta_{\mathcal{B}} - \theta_{\mathcal{H}})/\gamma\}$ through $\mathcal{B}$ labels.
  \item[(ii)] \emph{Optimal policy and tag value.} Welfare is $\max\{0, m\}$ without disclosure and $s_{\mathcal{B}}$ under an audit trail (before fixed costs). Under class tags, let $V_C(\gamma)$ denote the welfare gain over the uniform cut:
  \begin{itemize}
    \item If $m \ge 0$: $\theta_{\mathcal{B}} = 1$, label gap $\theta_{\mathcal{B}} - \theta_{\mathcal{H}} = \min\{1, \alpha_{\mathcal{H}} g\}$, and $V_C = \chi \alpha_{\mathcal{H}}^2 g / 2$ when $\alpha_{\mathcal{H}} g \le 1$, else $V_C = \chi \alpha_{\mathcal{H}} - \chi / (2g)$.
    \item If $m < 0$: $\theta_{\mathcal{H}} = 0$, $\theta_{\mathcal{B}} = \min\{1, s_{\mathcal{B}} g / \chi\}$, and $V_C = s_{\mathcal{B}}^2 g / (2\chi)$ when $s_{\mathcal{B}} g \le \chi$, else $V_C = s_{\mathcal{B}} - \chi / (2g)$.
  \end{itemize}
  In every case, the equilibrium hiding share lies strictly in $(0, 1)$, and at interior optima it is independent of $\gamma$: $z^* = \alpha_{\mathcal{H}}/(2\alpha_{\mathcal{H}} + b_{\mathcal{H}})$ when $m \ge 0$ and $z^* = s_{\mathcal{B}}/\{\chi(2\alpha_{\mathcal{H}} + b_{\mathcal{H}})\}$ when $m < 0$.
  \item[(iii)] \emph{Audit trail versus class tags.} $V_C(\gamma)$ is continuous and strictly increasing in $\gamma$, with $V_C(\gamma) \to 0$ as $\gamma \to 0$, $V_C(\gamma) \to V_A \equiv s_{\mathcal{B}} - \max\{0, m\}$ as $\gamma \to \infty$, and $V_C(\gamma) < V_A$ for all $\gamma > 0$. Given fixed costs $F_C \le F_A$, an owner-linked audit trail strictly dominates class tags if and only if $F_A - F_C < V_A - V_C(\gamma)$, i.e., whenever relabeling cost $\gamma$ lies below a unique threshold $\gamma^*$.
\end{enumerate}
\end{proposition}

The proof (Appendix~\ref{app:proofs}) yields three sharp economic takeaways. First, cheap relabeling restores observational equivalence in equilibrium: as $\gamma \to 0$, class tags become worthless. Second, the optimal targeted restriction is partial—when uniform retention is worthwhile ($m \ge 0$), the supervisor leaves the harmful label partly open; when $m < 0$, it also restricts the benign label because that label serves as the hiding channel. Third, an audit trail linking accounts to beneficial owners preserves its full value $V_A$ regardless of $\gamma$. While Proposition~\ref{prop:relabeling} adopts linear per-unit values $(n_{\mathcal{H}}, n_{\mathcal{B}})$ for closed-form transparency, Online Supplement Theorem~S15.4 proves that the same limits—collapse of class-tag value as $\gamma \to 0$ and convergence to $V_A$ as $\gamma \to \infty$—hold for arbitrary concave private payoffs, general convex concealment costs, and any Lipschitz social welfare functional $U(e)$, including quadratic aggregate-feedback stability externalities $U(e) = \sum_j \chi_j b_j e_j - \frac{\lambda}{2}(\sum_j \chi_j w_j e_j)^2$ directly linked to the feedback budget of Section~\ref{sec:equilibrium}.

Even without account relabeling, an untagged algorithm that anticipates targeting can disguise itself by shifting its trading rule toward its neighbors. When the identification divergence between the group and its neighbor scales at super-quadratic order $D^p$ ($p > 2$)—as in symmetric splitting ($p = 4$, Online Supplement Theorem~S2.3(ii)), two-group moment-matched hiding ($p = 6$, Theorem~\ref{thm:cost}), or a $K$-group moment-matched cluster ($p = 4K - 2$)—untagged surveillance becomes self-defeating:

\begin{theorem}[The surveillance paradox]\label{thm:paradox}
An untagged trader group chooses the distance $D \in [0, D_0]$ between its memory rate and its neighbor's within a moment-matched hiding family whose $N$-observation divergence scales as $J = c N D^p$ with exponent $p > 2$ (e.g., $p = 4$ for symmetric splitting, $p = 6$ for $K = 2$, and $p = 4K - 2$ for a $K$-group cluster). The group incurs strategy-distortion cost $k(D_0 - D)^2$ and targeting penalty $L > 0$ if detected by the supervisor's unbiased level-$\alpha$ test ($\pi(J) \ge \alpha$), with $D_0, k, L, c > 0$. Let $D^*$ minimize expected cost and $J^* = c N (D^*)^p$.
\begin{enumerate}
  \item[(i)] \emph{Incomplete hiding.} Against every unbiased test, $D^* > 0$, and $(D^*)^{p-2} \ge 2 k^2 D_0^2 / (2^p L^2 c N)$ whenever $D^* \le D_0 / 3$.
  \item[(ii)] \emph{Good tests force crowding.} If the test's power satisfies $\pi(J) - \alpha \ge \beta \min\{\sqrt{J}, \sqrt{\bar{J}}\}$ and stakes are high ($L \beta \sqrt{\bar{J}} > k D_0^2$), then for every $N$,
  \[
    (D^*)^{p-2} \le \frac{4 k^2 D_0^2}{L^2 \beta^2 c N}.
  \]
  \item[(iii)] \emph{More data, less information.} Under (ii), for all large $N$, $D^* \asymp N^{-1/(p-2)}$ and the supervisor's equilibrium information obeys $J^* \asymp N^{-2/(p-2)} \to 0$ (decaying as $N^{-1/2}$ for $p = 6$ and $N^{-1/(2K-2)}$ for $p = 4K - 2$). Detection probability converges to the false-alarm rate $\alpha$, while the group's strategy distortion approaches $k D_0^2$.
  \item[(iv)] \emph{Tags eliminate crowding.} Under a verified class tag, detection is independent of $D$, so $D^* = D_0$ and supervisory information $c N D_0^p$ grows linearly in $N$.
\end{enumerate}
\end{theorem}

The Neyman--Pearson test in a Gaussian experiment satisfies (ii) with $\beta = \sqrt{2}\phi(z_\alpha)$ and $\bar{J} = z_\alpha^2 / 2$ (proof in Appendix~\ref{app:proofs}). Without tags, expanding the price record provokes crowding faster than it accumulates statistical evidence, imposing a deadweight strategy distortion $k D_0^2$ while driving supervisory knowledge of composition to zero. By Lemma~\ref{lem:micro}(ii), the distortion cost $k(D_0 - D)^2$ arises directly as the leading-order profit loss of a Kalman-filtering desk deviating from its optimal memory rate $r_0$ in the small-group limit, while Theorem~\ref{thm:paradox} holds the neighbor's position fixed. When multiple untagged groups position strategically—including benign groups that may wish to separate—Online Supplement Proposition~S11.1(iv) shows that in any Nash equilibrium with bounded nearest-neighbor distances, every group's asymptotic detection probability remains bounded above by $\alpha + k_g \limsup_N \Delta_g^2 / L_g$; verified class tags eliminate both crowding and costly defensive repositioning simultaneously.

\subsection{Illustrative empirical calibration: Taiwan Stock Exchange tagged flows}

To illustrate the magnitude of the identification barrier in an actual exchange dataset, we calibrate Theorems~\ref{thm:cost} and~\ref{thm:forecast-kl} using daily class-tagged net buying from the Taiwan Stock Exchange (TWSE) across three reported institutional classes—foreign investors, investment trusts, and dealers—over 2013--2016 ($300$ stocks, $270{,}060$ stock-days). We estimate each class's persistence $r_c$ and response $w_c$ to the prior day's price deviation, along with the market's lagged-flow impact $\kappa$ and price persistence $\sum_{j=1}^5 a_j$, using stock fixed effects and stationary date-block bootstrap confidence intervals (Table~\ref{tab:taiwan}).

\begin{table}[!htbp]
\centering
\caption{Estimated tagged trading rules in Taiwan, 2013--2016 (300 stocks, 270,060 stock-days; stock fixed effects with 95\% date-block bootstrap intervals).}
\label{tab:taiwan}
\small
\begin{tabular}{@{}lccc@{}}
\toprule
Trader class & Persistence $r_c$ & Price response $w_c$ & Budget contribution $u_c = \kappa w_c / (1 - r_c)$ \\
\midrule
Foreign investors & $0.35 \; [0.31, 0.38]$ & $0.24 \; [0.21, 0.27]$ & $0.0014 \; [0.0012, 0.0016]$ \\
Investment trusts & $0.41 \; [0.39, 0.43]$ & $0.066 \; [0.057, 0.076]$ & $4.2 \times 10^{-4} \; [3.3 \times 10^{-4}, 5.3 \times 10^{-4}]$ \\
Dealers           & $0.33 \; [0.22, 0.42]$ & $-0.022 \; [-0.028, -0.017]$ & $-1.3 \times 10^{-4} \; [-1.8 \times 10^{-4}, -8.2 \times 10^{-5}]$ \\
\midrule
\multicolumn{4}{@{}l}{\footnotesize Lagged-flow price impact $\kappa = 0.0038 \; [0.0032, 0.0044]$; aggregate price persistence $\sum_{j=1}^5 a_j = 0.97$.} \\
\bottomrule
\end{tabular}
\end{table}

Granting a price-only supervisor full knowledge of the impact function $\kappa$ and the law of non-feedback order flow (so only class composition is unknown) yields three sharp quantitative contrasts (Online Supplement, Section~S19):
\begin{enumerate}
  \item \emph{Information gap.} Against the closest single-memory population, one stock-day of prices carries $1.7 \times 10^{-12}$~nats about class composition, whereas one stock-day of tagged flows carries $8.0 \times 10^{-7}$~nats—a factor of $4.7 \times 10^5$ more information. Reaching $80\%$ power at a $5\%$ test size from prices alone requires at least $2.7 \times 10^9$ stock-years ($9.1 \times 10^6$ years across 300 independent stocks).
  \item \emph{Targeted halt is unidentified from prices.} With tags, halting foreign feedback lowers the feedback budget by $0.0014$ ($95\%$ CI $[0.0012, 0.0016]$). Across all compositions statistically indistinguishable from the full panel price record (no test exceeding its size by more than $0.10$), the implied effect of halting one class spans $[-0.028, 0.028]$ (and $[-0.0027, 0.0027]$ even if each class is capped at twice the largest estimated contribution). Even with an infinite price record, permutation relabeling alone leaves $0.0014$, $4.2 \times 10^{-4}$, and $-1.3 \times 10^{-4}$ observationally equivalent.
  \item \emph{Uniform cut is pinned down by prices.} Over the exact same indistinguishable set, the effect of cutting all feedback uniformly lies within the tight band $[0.0014, 0.0019]$.
\end{enumerate}

Two caveats discipline this calibration. First, daily institutional flows exhibit small lagged price impact ($\kappa = 0.0038$), contributing $0.0017$ to an overall persistence budget of $0.97$ because daily aggregation blends intraday momentum with liquidity provision; thus Table~\ref{tab:taiwan} illustrates the \emph{relative} identification gap between uniform and targeted interventions rather than claiming that daily institutional flows destabilize Taiwanese equities. Second, as documented in Online Supplement Section~S6, a pre-registered out-of-sample reverse-causality diagnostic on 2019--2023 holdout data cannot rule out that daily price deviations partly reflect prior institutional order flow rather than pure causal feedback. Accordingly, these figures serve as an illustrative calibration of the relative identification gap between uniform and targeted cuts rather than a causal test of daily feedback destabilization; in intraday limit-order-book markets where aggregate feedback shares $\kappa \sum_i q_i$ are economically larger (as in the Section~\ref{sec:equilibrium} equilibrium calibrations underlying Figures~\ref{fig:near-instability} and~\ref{fig:intervention}), the same $(\Delta\eta)^{-(4K-2)}$ barrier operates with direct prudential significance.

\section{Scope and Conclusion}\label{sec:conclusion}

This paper characterizes the informational boundary between public market data and trader-level supervisory records in algorithmic feedback markets. Under both mechanical price impact and an endogenous inventory-dealer equilibrium, public prices reveal the population's aggregate feedback kernel—and nothing finer—at standard statistical rates, whereas separating closely spaced trader classes incurs a polynomial sample-complexity penalty of order $\Delta^{-(4K-2)}$ that persists even on the brink of instability.

Three policy conclusions follow. First, market-wide interventions (such as uniform feedback cuts or market-design changes calibrated to aggregate response) are statistically well-posed from public prices and are minimax-optimal when class identities are unverified. Second, targeted restrictions aimed at specific trader classes require class-tagged flow reports, which eliminate the spacing penalty for weight recovery. Third, because strategic traders can erode class tags by routing orders across labels or crowd their strategies together under untagged surveillance, durable targeted supervision requires owner-linked audit trails—such as the cross-venue customer linkage embedded in consolidated audit trail architectures.

Several limitations define natural directions for future work. Our explicit identification and upper-bound theorems rely on linear causal rules with finite-dimensional exponential or smooth memory; extensions to drifting populations, hard position limits, breach cascades, and saturating or learning algorithms are analyzed in the companion working papers (Niroula et al., 2026b, 2026c). Our passive sample-complexity bounds assume Gaussian (or quadratically regular) innovations and exclude bounded-support noise laws. Finally, while our contribution is theoretical and calibrated to public exchange data, a direct causal empirical test of the supervisory separation would evaluate an exogenous, externally documented class-specific trading restriction against pre-specified predictions.

\section*{Declarations}

\textbf{Author contributions.} Pratik Niroula: conceptualization, methodology, formal analysis, software, investigation, data curation, writing (original draft). Prashant Niroula: validation (proofs, Lean~4 formalizations, and empirical calculations), writing (review and editing). Sujan Khatiwada: validation (proofs, Lean~4 formalizations, and empirical calculations), writing (review and editing). All authors approved the final manuscript.

\textbf{Funding and competing interests.} This research received no external funding. The authors declare no competing interests.

\textbf{Data and code availability.} All empirical data are publicly available from the Taiwan Stock Exchange (\url{https://www.twse.com.tw}). Replication scripts, numerical checks, and Lean~4 formalization files are archived at \url{https://github.com/Pratikn03/Fianince}.

\textbf{Declaration of generative AI and AI-assisted technologies.} During the preparation of this manuscript, the authors used large language model tools (Claude, ChatGPT/Codex, and Gemini) to assist with literature cross-checking, numerical verification scripts, Lean~4 formalization coding, and language editing. All mathematical proofs, economic arguments, empirical calculations, and formal statements were independently verified and edited by the authors, who take full responsibility for the content of the publication.

\appendix
\section{Proofs of Main Economic Results}\label{app:proofs}

\begin{proof}[Proof of Theorem~\ref{thm:sufficiency}]
We argue by strong induction on $t \le T$ for any fixed strategy profile $(X, P)$ and realization $(v, z)$. Suppose prices and total order flows coincide across the two economies for all dates $j < t$. Then: (1) public signals $s_j = \mathcal{S}_j(p_0, \dots, p_j)$ coincide for all $j < t$; (2) aggregate feedback flow $f_t = \sum_{m < t} k_m s_{t-1-m}$ coincides because $k_m = k'_m$ for $m < T$; (3) the insider's order $X_t(v, \mathcal{H}_{t-1})$ coincides by causality; and (4) total order flow $Y_t = X_t + z_t + f_t$ and the dealer's price $P_t(\mathcal{H}_{t-1}, Y_t)$ coincide by causality. Hence price and flow paths, realized insider profit $\sum_{t < T}(v - p_t) X_t$, and dealer inventory $-\sum_{t < T} Y_t$ coincide pathwise. Because strategy sets and expected payoffs over $(v, z)$ are identical, best-response correspondences and equilibria coincide under any solution concept defined from public histories and payoffs. For any finite linear population with individual kernels $\kappa^{(j)}$, $f_t = \sum_j \sum_{m < t} \kappa_m^{(j)} s_{t-1-m} = \sum_{m < t} k_m s_{t-1-m}$ with $k = \sum_j \kappa^{(j)}$, so the identical induction applies.
\end{proof}

\begin{proof}[Proof of Theorem~\ref{thm:inv-eq}]
\emph{(i)} Conjecture $x_t = -\kappa H_t + \sum_i \beta_i u_{i,t+1}$ and set $D = 1 - \sum_i \beta_i w_i$. Using $\E_t[\xi_{t+1}] = 0$, $\E_t[H_{t+1}] = \bar{a} H_t - \sum_i u_{i,t+1}$, and $\E_t[u_{i,t+2}] = r_i u_{i,t+1} + w_i \E_t[x_{t+1}]$, taking conditional expectations gives
\[
  D\, \E_t[x_{t+1}] = -\kappa \bar{a} H_t + \sum_i (\kappa + \beta_i r_i) u_{i,t+1}.
\]
The dealer optimality condition $x_t = \E_t[x_{t+1}] - A V H_t$ must hold identically in $(H_t, \{u_{i,t+1}\})$. Matching coefficients on $u_{i,t+1}$ yields $\beta_i = (\kappa + \beta_i r_i)/D$, so $\beta_i = \kappa / (D - r_i)$. Substituting into $D = 1 - \sum_i \beta_i w_i$ gives $D = 1 - \kappa G(D)$, and matching coefficients on $H_t$ gives $\kappa(D - \bar{a}) = A V D$. For the one-step innovation $\delta = x_{t+1} - \E_t[x_{t+1}]$, we have $\delta = \kappa \xi_{t+1} + \sum_i \beta_i w_i \delta$, hence $\delta = (\kappa/D)\xi_{t+1}$ and $V = \sigma_F^2 + \kappa^2 \sigma_\xi^2 / D^2$. Conversely, any stable solution of~\eqref{eq:eq-system} satisfies dealer optimality and variance consistency in every state.

\emph{(ii)} Setting $w_i = 0$ yields $G \equiv 0$, $D = 1$, $x_t = -\kappa H_t = \bar{a} x_{t-1} + \kappa Y_t$, and the quadratic $A \sigma_\xi^2 \kappa^2 - (1 - \bar{a})\kappa + A \sigma_F^2 = 0$, whose discriminant is $\Lambda_0 = (1 - \bar{a})^2 - 4 A^2 \sigma_F^2 \sigma_\xi^2$.

\emph{(iii)} System~\eqref{eq:eq-system} depends on the population solely through $G(D) = \int (D - r)^{-1} dM(r)$, hence solely through $M$.

\emph{(iv)} Define $\Phi(\kappa, D; M) = \bigl(D + \kappa G(D) - 1,\; \kappa(D - \bar{a}) - A(\sigma_F^2 D + \kappa^2 \sigma_\xi^2 / D)\bigr)$. At $M = 0$, $\Phi(\kappa_-, 1; 0) = 0$ with triangular Jacobian $\partial_D \Phi_1 = 1$, $\partial_\kappa \Phi_1 = 0$, and $\partial_\kappa \Phi_2 = (1 - \bar{a}) - 2 A \kappa_- \sigma_\xi^2 = \sqrt{\Lambda_0} > 0$. The implicit function theorem yields a unique continuous branch near $(\kappa_-, 1)$ for small $G$. Linearizing gives $dD = -\kappa_- G(1)$ and $d\kappa = -\kappa_-(\bar{a} + 2 A \kappa_- \sigma_\xi^2)\, dD / \sqrt{\Lambda_0}$, with $G(1) = \sum_i w_i/(1 - r_i) = \sum_i q_i$. Stability holds by continuity of eigenvalues from $M = 0$.
\end{proof}

\begin{proof}[Proof of Theorem~\ref{thm:eq-id}]
\emph{(i)} Write $D x_t = -\kappa H_t + \sum_i \beta_i r_i u_{i,t}$ and $x_{t-1} = -\kappa H_{t-1} + \sum_i \beta_i u_{i,t}$, where $\beta_i = \kappa/(D - r_i)$. Substituting $H_t = \bar{a} H_{t-1} - \sum_i u_{i,t} - \xi_t$ yields
\[
  D x_t - \bar{a} x_{t-1} = \kappa \xi_t + \sum_i \bigl(\kappa + \beta_i(r_i - \bar{a})\bigr) u_{i,t} = \kappa \xi_t + \sum_i \kappa \frac{D - \bar{a}}{D - r_i}\, u_{i,t}.
\]
Dividing by $D$ and defining $\tilde{u}_{i,t} = \frac{D - \bar{a}}{D - r_i} u_{i,t}$ gives $x_t = a^* x_{t-1} + \kappa^*(\xi_t + \sum_i \tilde{u}_{i,t})$ and $\tilde{u}_{i,t+1} = r_i \tilde{u}_{i,t} + \tilde{w}_i x_t$ pathwise, with $(a^*, \kappa^*, \tilde{w}_i)$ as stated. Because $\kappa \neq 0$ and $D \neq \bar{a}$, the state transformation is invertible and preserves stability.

\emph{(ii)}-- \emph{(iii)} Part~(ii) follows from Online Supplement Theorem~S4.3(i) applied to the closed loop in~(i). For~(iii), $G(D) = \sum_i \tilde{w}_i/(D - \bar{a})$ turns $D = 1 - \kappa G(D)$ into $D = 1 - \kappa^* \tilde{M}([0,1])/(1 - a^*)$, which recovers $D$, $\bar{a} = a^* D$, $\kappa = \kappa^* D$, and $M = h^{-1} \tilde{M}$ uniquely.

\emph{(iv)} Given any hard pair $(a^*, \kappa^*, W, r_c, \Delta\eta)$ satisfying the stability condition $a^* + \kappa^* W / (1 - r_c - \Delta\eta) < 1$, we have $\kappa^* W / (1 - a^*) < 1 - r_c - \Delta\eta$, so setting $D = 1 - \kappa^* W/(1 - a^*)$ ensures $D > r_c + \Delta\eta \ge \max_i r_i$. Choosing $\bar{a} = a^* D$, $\kappa = \kappa^* D$, any $\sigma_F^2, \sigma_\xi^2 > 0$, $A = \kappa^* D(1 - a^*)/(\sigma_F^2 + \kappa^{*2} \sigma_\xi^2)$, and $w_i = \tilde{w}_i(D - r_i)/(D - \bar{a}) > 0$ realizes both economies as equilibria with identical primitives $(A, \sigma_F^2, \sigma_\xi^2, \bar{a})$ and identical $(\kappa, D)$, preserving the $(\Delta\eta)^{-(4K-2)}$ lower bound.
\end{proof}

\begin{proof}[Proof of Theorem~\ref{thm:smooth}]
Let $G_K$ denote the $K$-point Gauss--Legendre quadrature rule on $[-1, 1]$ normalized to unit mass, and let $R_K$ denote the $K$-point Gauss--Radau rule with fixed endpoint $-1$, also normalized to unit mass. Both measures have positive weights; $G_K$ integrates every polynomial of degree at most $2K - 1$ exactly, and $R_K$ integrates every polynomial of degree at most $2K - 2$ exactly. Hence $\int P\, dG_K = \int P\, dR_K$ for every polynomial $P$ of degree at most $2K - 2$. Because $R_K$ places positive mass at $-1$ while $G_K$ does not, their 1-Wasserstein distance satisfies $d_K = W_1(G_K, R_K) > 0$. Define $\nu$ and $\nu'$ as the pushforwards of $W G_K$ and $W R_K$ under $s \mapsto \theta_c + \Delta s$; then both have $K$ atoms in $[\theta_c - \Delta, \theta_c + \Delta]$, total mass $W$, identical moments of orders $0, \dots, 2K - 2$ about $\theta_c$, and $W_1(\nu, \nu') = d_K W \Delta$.

\emph{(i)} Fix any history $h$. By Taylor's theorem with Lagrange remainder, $f_\theta(h) = P_h(\theta) + \rho_h(\theta)$, where $P_h$ is a polynomial of degree $2K - 2$ in $\theta - \theta_c$ and $|\rho_h(\theta)| \le B(h)|\theta - \theta_c|^{2K-1}/(2K - 1)!$. Since $P_h$ integrates identically under $\nu$ and $\nu'$,
\[
  |F_\nu(h) - F_{\nu'}(h)|
  = \Bigl|\int \rho_h\, d\nu - \int \rho_h\, d\nu'\Bigr|
  \le \int |\rho_h|\, d\nu + \int |\rho_h|\, d\nu'
  \le \frac{2W}{(2K - 1)!}\, B(h)\, \Delta^{2K-1}.
\]

\emph{(ii)} Let $P$ and $P'$ be the laws of the $T$-period record under $\nu$ and $\nu'$ from the common initial law and common predictable probe policy. Conditional on the past and probe $v_t$, $x_t$ is Gaussian with variance $\sigma^2$ and conditional means differing by $\kappa(F_\nu(h_{t-1}) - F_{\nu'}(h_{t-1}))$. By the chain rule for Kullback--Leibler divergence, part~(i), and the uniform moment bound $\E_\nu[B(h_t)^2] \le \bar{B}^2$,
\[
  \KL(P \| P')
  = \sum_{t=1}^T \E_P\Bigl[\frac{\kappa^2 \bigl(F_\nu(h_{t-1}) - F_{\nu'}(h_{t-1})\bigr)^2}{2\sigma^2}\Bigr]
  \le T\, \frac{2 \kappa^2 W^2 \bar{B}^2}{\sigma^2 ((2K - 1)!)^2}\, \Delta^{4K-2}.
\]
Any test with both error probabilities strictly below $1/4$ satisfies $\TV(P, P') > 1/2$, which by Pinsker's inequality implies $\KL(P \| P') > 1/2$. Rearranging yields the stated lower bound on $T$.
\end{proof}

\begin{proof}[Proof of Theorem~\ref{thm:forecast-kl}]
Conditional on the observed history $h_t$ and predictable probe $v_t$, $x_{t+1}$ has variance $\sigma^2$ under both $P$ and $P'$, with conditional means differing by $d_t = \kappa(F(h_t) - F'(h_t))$. The one-step Gaussian KL divergence is $d_t^2 / (2\sigma^2)$. Since initial laws and conditional probe laws match, the chain rule gives $\KL(P_T \| P'_T) = (2\sigma^2)^{-1} \sum_{t < T} \E_P[d_t^2] = \Lambda_T / 2$. Under $P$, the forecast error of the alternative conditional mean is $\varepsilon_{t+1} + d_t$, whose conditional mean-squared error is $\sigma^2 + d_t^2$, so excess MSE is $V_t = \E_P[d_t^2]$. If both testing errors are strictly below $1/4$, their sum is below $1/2$, so $\TV(P_T, P'_T) > 1/2$, and Pinsker's inequality gives $\KL(P_T \| P'_T) > 1/2$, i.e., $\Lambda_T > 1$. Finally, $q(X_{1:T}) = \sum_j \pi_j p_j(X_{1:T}) \ge \pi_* p_*(X_{1:T})$ pointwise, yielding $\log(p_*/q) \le \log(1/\pi_*)$.
\end{proof}

\begin{proof}[Proof of Theorems~\ref{thm:minimax} and~\ref{thm:partial-minimax} and Corollary~\ref{cor:halt}]
For Theorem~\ref{thm:minimax}, averaging over all $K!$ permutations $\sigma \in S_K$ places every label on every class equally often:
\[
  \frac{1}{K!} \sum_{\sigma \in S_K} S'_\sigma = a + \sum_{c=1}^K \theta_c \Bigl(\frac{1}{K} \sum_{d=1}^K u_d\Bigr) = a + \bar{\theta} U,
  \qquad
  \frac{1}{K!} \sum_{\sigma \in S_K} \theta \circ \sigma = \bar{\theta}\mathbf{1}.
\]
By convexity of $\ell$ and $C$ (and permutation invariance of $C$), Jensen's inequality gives part~(i):
\[
  \ell(a + \bar{\theta} U) + C(\bar{\theta}\mathbf{1})
  \le \frac{1}{K!} \sum_{\sigma} \ell(S'_\sigma) + C(\theta)
  \le \max_{\sigma} \ell(S'_\sigma) + C(\theta).
\]
For part~(ii), suppose $\theta \in [0, 1]^K$ is non-uniform, so $\theta_{c_1} \neq \theta_{c_2}$ for some $c_1 \neq c_2$. Then the transposition $\tau = (c_1\, c_2)$ satisfies $\theta \circ \tau \neq \theta$. Since $\bar{\theta}\mathbf{1} = \frac{1}{K!}\sum_{\sigma \in S_K} (\theta \circ \sigma)$ is a strictly positive convex combination of points that are not all equal, strict convexity and permutation invariance of $C$ yield $C(\bar{\theta}\mathbf{1}) < \frac{1}{K!}\sum_\sigma C(\theta \circ \sigma) = C(\theta)$, giving strict dominance. (Similarly, if $u_{d_1} \neq u_{d_2}$ and $\theta_{c_1} \neq \theta_{c_2}$, swapping the assignment of $d_1, d_2$ to $c_1, c_2$ produces distinct budgets $S'_{\sigma_1} \neq S'_{\sigma_2}$, so strict convexity of $\ell$ yields $\ell(a + \bar{\theta} U) < \frac{1}{K!}\sum_\sigma \ell(S'_\sigma) \le \max_\sigma \ell(S'_\sigma)$.) Hence every minimax-optimal policy must be uniform; on the compact interval $[0, 1]$, the strictly convex continuous function $\bar{\theta} \mapsto \ell(a + \bar{\theta} U) + C(\bar{\theta}\mathbf{1})$ attains a unique minimizer $\bar{\theta}^*$. Thus $\bar{\theta}^*\mathbf{1}$ is the unique minimax supervisory policy.
For Corollary~\ref{cor:halt}, halting label $F$ yields post-halt budget $S - u_{\sigma(F)}$, whose worst case over $\sigma$ is $S - \min_d u_d$ (attained whenever $\sigma(F) \in \arg\min_d u_d$). A uniform cut $\theta = (\bar{s} - a)/(S - a)$ yields budget $a + \theta(S - a) = \bar{s}$ exactly and removes total feedback $(1 - \theta)(S - a) = S - \bar{s}$, the algebraic minimum needed to reach $\bar{s}$. For Theorem~\ref{thm:partial-minimax}, restrict the averaging group to the product of within-group permutations $\prod_{g=1}^G S_g$: within each group $g$, every label averages to $\bar{\theta}_g$, so the averaged budget is $a + \sum_g \bar{\theta}_g U_g$ and the identical Jensen argument yields the minimax bound and $V_{\mathrm{tag}} \ge 0$.
\end{proof}

\begin{proof}[Proof of Lemma~\ref{lem:micro}]
\emph{(i)} Linearizing social welfare around the no-feedback equilibrium of Theorem~\ref{thm:inv-eq}(iv), each retained unit of class $j$ yields private surplus $b_j$ and raises aggregate feedback intensity by $\phi_j$, increasing equilibrium impact $\kappa$ proportionally to $\kappa_-^2(\bar{a} + 2A\kappa_-\sigma_\xi^2)/\sqrt{\Lambda_0} > 0$. Valuing the marginal market-wide impact and inventory-risk externality at $\lambda_{\mathrm{ext}} > 0$ gives the first-order net social value $n_j = b_j - \lambda_{\mathrm{ext}}\phi_j$, with $n_{\mathcal{H}} < 0 < n_{\mathcal{B}}$ under the stated sign conditions.

\emph{(ii)} Under $\hat{v}_{t+1} = r \hat{v}_t + (1 - r)(v_t + \nu_t)$, the forecast error $e_t = v_t - \hat{v}_t$ evolves as $e_{t+1} = r e_t + \zeta_{t+1} - (1 - r)\nu_t$. Taking stationary variances with independent innovations gives $P(r) = r^2 P(r) + \sigma_\zeta^2 + (1 - r)^2 \sigma_\nu^2$, so
\[
  P(r) = \frac{\sigma_\zeta^2 + (1 - r)^2 \sigma_\nu^2}{1 - r^2}
       = \frac{\sigma_\zeta^2}{1 - r^2} + \frac{2\sigma_\nu^2}{1 + r} - \sigma_\nu^2,
\]
using $(1 - r)^2 / (1 - r^2) = (1 - r)/(1 + r) = 2/(1 + r) - 1$. Differentiating twice on $(0, 1)$ gives
\[
  P'(r) = \frac{2 r \sigma_\zeta^2}{(1 - r^2)^2} - \frac{2\sigma_\nu^2}{(1 + r)^2},
  \qquad
  P''(r) = \frac{2\sigma_\zeta^2(1 + 3r^2)}{(1 - r^2)^3} + \frac{4\sigma_\nu^2}{(1 + r)^3} > 0.
\]
Since $P'(0) = -2\sigma_\nu^2 < 0$, $\lim_{r \uparrow 1} P'(r) = +\infty$, and $P''(r) > 0$ throughout $(0, 1)$, $P$ is strictly convex with a unique interior minimizer $r_0 \in (0, 1)$ satisfying $P'(r_0) = 0$. Substituting the one-period optimal trade $u_t = \hat{v}_t / (\kappa + \lambda)$ into expected profit gives $\Pi(r) = \frac{1}{2(\kappa + \lambda)}\bigl(\E[v_t^2] - P(r)\bigr)$ (since $\E[v_t \hat{v}_t] - \frac{1}{2}\E[\hat{v}_t^2] = \frac{1}{2}\E[v_t^2] - \frac{1}{2}\E[(v_t - \hat{v}_t)^2]$), so
\[
  \Pi(r_0) - \Pi(r) = \frac{P(r) - P(r_0)}{2(\kappa + \lambda)}
  = \frac{P''(r_0)}{4(\kappa + \lambda)}(r - r_0)^2 + O\!\left(|r - r_0|^3\right).
\]
Finally, by Theorem~\ref{thm:inv-eq}(iv), $\partial \kappa / \partial r_i = O(w_i)$, so for a small group ($w_i \ll 1$) the induced equilibrium shift in $\kappa$ is of order $O(w_i)$, leaving $k = P''(r_0)/(4(\kappa + \lambda))$ invariant at leading order.
\end{proof}

\begin{proof}[Proof of Proposition~\ref{prop:relabeling}]
\emph{(i)} An $\mathcal{H}$ trader maximizes $b_{\mathcal{H}}\{(1 - z)\theta_{\mathcal{H}} + z \theta_{\mathcal{B}}\} - \gamma z^2/2 = b_{\mathcal{H}}\theta_{\mathcal{H}} + b_{\mathcal{H}}(\theta_{\mathcal{B}} - \theta_{\mathcal{H}})z - \gamma z^2/2$ over $z \in [0, 1]$, which is strictly concave with unique maximizer $z^* = \min\{1,\, b_{\mathcal{H}}(\theta_{\mathcal{B}} - \theta_{\mathcal{H}})/\gamma\}$.

\emph{(ii)} Without disclosure, welfare is $m\theta$, maximized at $\max\{0, m\}$. Under an audit trail ($z \equiv 0$), welfare $\chi \theta_{\mathcal{H}} n_{\mathcal{H}} + s_{\mathcal{B}} \theta_{\mathcal{B}}$ is maximized at $(\theta_{\mathcal{H}}, \theta_{\mathcal{B}}) = (0, 1)$ with value $s_{\mathcal{B}}$. Under class tags, write $d = \theta_{\mathcal{B}} - \theta_{\mathcal{H}} \in [0, \theta_{\mathcal{B}}]$. Any policy with $b_{\mathcal{H}} d \ge \gamma$ induces $z^* = 1$ and welfare $m\theta_{\mathcal{B}} - \chi\gamma/2 < m\theta_{\mathcal{B}}$, strictly dominated by $d = 0$. For $b_{\mathcal{H}} d < \gamma$, substituting $z^* = b_{\mathcal{H}} d / \gamma$ into social welfare $m\theta_{\mathcal{B}} + \chi \alpha_{\mathcal{H}} d (1 - z^*) - \chi \gamma (z^*)^2 / 2$ yields
\[
  W(\theta_{\mathcal{B}}, d) = m\theta_{\mathcal{B}} + \chi \Bigl(\alpha_{\mathcal{H}} d - \frac{d^2}{2g}\Bigr),
  \qquad g = \frac{\gamma}{b_{\mathcal{H}}(2\alpha_{\mathcal{H}} + b_{\mathcal{H}})}.
\]
Maximizing over $d \in [0, \theta_{\mathcal{B}}]$ gives $d^* = \min\{\theta_{\mathcal{B}}, \hat{d}\}$ where $\hat{d} = \alpha_{\mathcal{H}} g < \gamma / b_{\mathcal{H}}$ (so $b_{\mathcal{H}} d < \gamma$ never binds). The concentrated objective $\varphi(\theta_{\mathcal{B}}) = s_{\mathcal{B}}\theta_{\mathcal{B}} - \chi \theta_{\mathcal{B}}^2 / (2g)$ for $\theta_{\mathcal{B}} \le \hat{d}$ and $m\theta_{\mathcal{B}} + \chi \alpha_{\mathcal{H}}^2 g / 2$ for $\theta_{\mathcal{B}} \ge \hat{d}$ is concave and $C^1$. If $m \ge 0$, $\varphi$ is nondecreasing, so $\theta_{\mathcal{B}} = 1$ and $d^* = \min\{1, \alpha_{\mathcal{H}} g\}$. If $m < 0$, $\varphi'(\theta_{\mathcal{B}}) = 0$ at $\hat{\theta}_{\mathcal{B}} = s_{\mathcal{B}} g / \chi < \hat{d}$, so $\theta_{\mathcal{B}} = \min\{1, s_{\mathcal{B}} g / \chi\}$ and $\theta_{\mathcal{H}} = 0$. Substituting $d^*$ into $z^* = b_{\mathcal{H}} d^* / \gamma$ gives $z^* \in (0, \alpha_{\mathcal{H}}/(2\alpha_{\mathcal{H}} + b_{\mathcal{H}})] \subset (0, 1)$, independent of $\gamma$ at interior optima.

\emph{(iii)} Each branch of $V_C(\gamma)$ is strictly increasing in $g \propto \gamma$, continuous at the kink ($\alpha_{\mathcal{H}} g = 1$ or $s_{\mathcal{B}} g = \chi$), satisfies $V_C(\gamma) \to 0$ as $\gamma \to 0$ and $V_C(\gamma) \uparrow V_A$ as $\gamma \to \infty$, and remains strictly below $V_A$ for all finite $\gamma > 0$. Hence $V_A - V_C(\gamma)$ decreases strictly from $V_A$ to $0$, yielding the threshold characterization.
\end{proof}

\begin{proof}[Proof of Theorem~\ref{thm:paradox}]
\emph{(i)} Moving from $D = 0$ to small $D > 0$ changes expected cost by at most $-2k D_0 D + k D^2 + L \sqrt{c N D^p / 2}$ (since $\pi(J) - \alpha \le \TV \le \sqrt{J/2}$ by Pinsker's inequality). Because $p > 2$, the linear term $-2k D_0 D$ dominates near zero, so $D^* > 0$. If $D^* \le D_0 / 3$, then $2D^* \le D_0$ is feasible; optimality of $D^*$ against $2D^*$, together with $\pi(J^*) \ge \alpha$ and Pinsker's inequality, gives
\[
  k D_0 D^* \le k\bigl[(D_0 - D^*)^2 - (D_0 - 2D^*)^2\bigr] \le L\sqrt{c N (2D^*)^p / 2}.
\]
Squaring and dividing by $(D^*)^2$ yields $(D^*)^{p-2} \ge 2 k^2 D_0^2 / (2^p L^2 c N)$.

\emph{(ii)} Optimality of $D^*$ against $D = 0$ (where $J = 0$ and $\pi(0) = \alpha$) gives
\[
  L(\pi(J^*) - \alpha) \le k\bigl[D_0^2 - (D_0 - D^*)^2\bigr] \le \min\{k D_0^2,\; 2 k D_0 D^*\}.
\]
If $J^* \ge \bar{J}$, the power envelope implies $L \beta \sqrt{\bar{J}} \le k D_0^2$, contradicting the high-stakes condition. Thus $J^* < \bar{J}$ and $L \beta \sqrt{c N (D^*)^p} \le 2 k D_0 D^*$; squaring and dividing by $(D^*)^2$ gives $(D^*)^{p-2} \le 4 k^2 D_0^2 / (L^2 \beta^2 c N)$.

\emph{(iii)} For all sufficiently large $N$, part~(ii) forces $D^* \le D_0 / 3$, so both~(i) and~(ii) hold simultaneously: $(D^*)^{p-2} \asymp N^{-1}$, i.e., $D^* \asymp N^{-1/(p-2)}$, and $J^* = c N (D^*)^p \asymp N \cdot N^{-p/(p-2)} = N^{-2/(p-2)} \to 0$. Hence $\pi(J^*) \le \alpha + \sqrt{J^*/2} \to \alpha$ and $k(D_0 - D^*)^2 \to k D_0^2$.

\emph{(iv)} Under a class tag, $\pi$ does not depend on $D$, so $k(D_0 - D)^2$ is uniquely minimized at $D^* = D_0$ and $J^* = c N D_0^p$.
\end{proof}

\section*{References}
\begin{list}{}{\small\setlength{\leftmargin}{1.8em}\setlength{\itemindent}{-1.8em}\setlength{\itemsep}{3pt}\setlength{\parsep}{0pt}}

\item Auclert, A., Rognlie, M., Straub, L., 2024. The intertemporal Keynesian cross. \emph{Journal of Political Economy} 132(12), 4068--4121.

\item Batenkov, D., Goldman, G., Yomdin, Y., 2021. Super-resolution of near-colliding point sources. \emph{Information and Inference: A Journal of the IMA} 10(2), 515--572.

\item Blackwell, D., 1953. Equivalent comparisons of experiments. \emph{Annals of Mathematical Statistics} 24(2), 265--272.

\item Boswijk, H.P., Hommes, C.H., Manzan, S., 2007. Behavioral heterogeneity in stock prices. \emph{Journal of Economic Dynamics and Control} 31(6), 1938--1970.

\item Brock, W.A., Hommes, C.H., 1998. Heterogeneous beliefs and routes to chaos in a simple asset pricing model. \emph{Journal of Economic Dynamics and Control} 22(8--9), 1235--1274.

\item Budish, E., Cramton, P., Shim, J., 2015. The high-frequency trading arms race: Frequent batch auctions as a market design response. \emph{Quarterly Journal of Economics} 130(4), 1547--1621.

\item Cagan, P., 1956. The monetary dynamics of hyperinflation. In: Friedman, M.\ (Ed.), \emph{Studies in the Quantity Theory of Money}. University of Chicago Press, pp.\ 25--117.

\item Celov, D., Leipus, R., Philippe, A., 2010. Asymptotic normality of the mixture density estimator in a disaggregation scheme. \emph{Journal of Nonparametric Statistics} 22(4), 425--442.

\item Crocker, K.J., Morgan, J., 1998. Is honesty the best policy? Curtailing insurance fraud through optimal incentive contracts. \emph{Journal of Political Economy} 106(2), 355--375.

\item Dawid, A.P., 1984. Statistical theory: The prequential approach. \emph{Journal of the Royal Statistical Society, Series A} 147(2), 278--292.

\item Donoho, D.L., 1994. Statistical estimation and optimal recovery. \emph{Annals of Statistics} 22(1), 238--270.

\item Forni, M., Lippi, M., 1997. \emph{Aggregation and the Microfoundations of Dynamic Macroeconomics}. Oxford University Press.

\item Franke, R., Westerhoff, F., 2012. Structural stochastic volatility in asset pricing dynamics: Estimation and model contest. \emph{Journal of Economic Dynamics and Control} 36(8), 1193--1211.

\item Frino, A., Johnstone, D., Zheng, H., 2010. Anonymity, stealth trading, and the information content of broker identity. \emph{Financial Review} 45(3), 501--522.

\item Granger, C.W.J., 1980. Long memory relationships and the aggregation of dynamic models. \emph{Journal of Econometrics} 14(2), 227--238.

\item Heinrich, P., Kahn, J., 2018. Strong identifiability and optimal minimax rates for finite mixture estimation. \emph{Annals of Statistics} 46(6A), 2844--2870.

\item Hendry, D.F., Hubrich, K., 2011. Combining disaggregate forecasts or combining disaggregate information to forecast an aggregate. \emph{Journal of Business \& Economic Statistics} 29(2), 216--227.

\item Hommes, C.H., 2006. Heterogeneous agent models in economics and finance. In: Tesfatsion, L., Judd, K.L.\ (Eds.), \emph{Handbook of Computational Economics}, Vol.~2. Elsevier, pp.\ 1109--1186.

\item Ing, C.-K., Wei, C.-Z., 2005. Order selection for same-realization predictions in autoregressive processes. \emph{Annals of Statistics} 33(5), 2423--2474.

\item Kirilenko, A., Kyle, A.S., Samadi, M., Tuzun, T., 2017. The Flash Crash: High-frequency trading in an electronic market. \emph{Journal of Finance} 72(3), 967--998.

\item Koyck, L.M., 1954. \emph{Distributed Lags and Investment Analysis}. North-Holland, Amsterdam.

\item Lacker, J.M., Weinberg, J.A., 1989. Optimal contracts under costly state falsification. \emph{Journal of Political Economy} 97(6), 1345--1363.

\item Le Cam, L., 1964. Sufficiency and approximate sufficiency. \emph{Annals of Mathematical Statistics} 35(4), 1419--1455.

\item Linnainmaa, J.T., Saar, G., 2012. Lack of anonymity and the inference from order flow. \emph{Review of Financial Studies} 25(5), 1414--1456.

\item L\"utkepohl, H., 1987. \emph{Forecasting Aggregated Vector ARMA Processes}. Lecture Notes in Economics and Mathematical Systems 284. Springer, Berlin.

\item Lux, T., Zwinkels, R.C.J., 2018. Empirical validation of agent-based models. In: Hommes, C., LeBaron, B.\ (Eds.), \emph{Handbook of Computational Economics}, Vol.~4. Elsevier, pp.\ 437--488.

\item Menkveld, A.J., 2016. The economics of high-frequency trading: Taking stock. \emph{Annual Review of Financial Economics} 8, 1--24.

\item Muth, J.F., 1960. Optimal properties of exponentially weighted forecasts. \emph{Journal of the American Statistical Association} 55(290), 299--306.

\item Nerlove, M., 1958. Adaptive expectations and cobweb phenomena. \emph{Quarterly Journal of Economics} 72(2), 227--240.

\item Niroula, P., Alkhushayni, S., Chen, Y., Bhavsar, M., Niroula, P., 2026a. From human markets to machine markets: Transient amplification, an exact depth condition, and a maximal pre-trade guard. Working paper, Minnesota State University, Mankato.

\item Niroula, P., Niroula, P., Khatiwada, S., 2026b. After a breach, and beyond linear rules: Companion results to ``Observationally equivalent algorithms''. Working paper, \url{https://github.com/Pratikn03/Fianince}.

\item Niroula, P., Niroula, P., Khatiwada, S., 2026c. Observationally equivalent algorithms: Full working paper with all supplementary results. Working paper, \url{https://github.com/Pratikn03/Fianince}.

\item Pesaran, M.H., Chudik, A., 2014. Aggregation in large dynamic panels. \emph{Journal of Econometrics} 178(2), 273--285.

\item Sanders, D.R., Irwin, S.H., Merrin, R.P., 2009. Smart money: The forecasting ability of CFTC large traders in agricultural futures markets. \emph{Journal of Agricultural and Resource Economics} 34(2), 276--296.

\item Sargent, T.J., 1976. The observational equivalence of natural and unnatural rate theories of macroeconomics. \emph{Journal of Political Economy} 84(3), 631--640.

\item Securities and Exchange Commission, 2012. Consolidated Audit Trail. Final Rule, Release No.\ 34-67457, 18 July 2012.

\item van Waarde, H.J., Eising, J., Trentelman, H.L., Camlibel, M.K., 2020. Data informativity: A new perspective on data-driven analysis and control. \emph{IEEE Transactions on Automatic Control} 65(11), 4753--4768.

\item Vovk, V., 2001. Competitive on-line statistics. \emph{International Statistical Review} 69(2), 213--248.

\item Wu, Y., Yang, P., 2020. Optimal estimation of Gaussian mixtures via denoised method of moments. \emph{Annals of Statistics} 48(4), 1981--2007.

\item Zaffaroni, P., 2004. Contemporaneous aggregation of linear dynamic models in large economies. \emph{Journal of Econometrics} 120(1), 75--102.

\end{list}

\end{document}
